% 本文件由 LaTeX 排版 Agent 根据有效 translation.json 自动生成。
% author=Ambrose; work_id=3403; title=论撒提鲁斯之死

\chapter{论\Person{撒提鲁斯}之死}

% structure: p0001 source=h1 target=omitted reason=page-title-already-rendered-by-parent

第一卷\newline{}第二卷\par

\SourceRecord{Translated by H. de Romestin, E. de Romestin and H.T.F. Duckworth.；Nicene and Post-Nicene Fathers, Second Series；Vol. 10.；Edited by Philip Schaff and Henry Wace.；Buffalo, NY: Christian Literature Publishing Co.,；1896.；<http://www.newadvent.org/fathers/3403.htm>.}

% structure: p0001 source=h1 target=section reason=page-title

\section{论撒提若之死（卷一）}

% structure: p0002 source=h2 target=subsection reason=relative-heading-rank; base=h2

\subsection{引言}

除他的姐姐\Person{玛尔策利纳}于圣诞节（约公元353年）从教宗\Person{利贝留}手中领受面纱外，\Person{圣盎博罗削}还有一个弟弟名叫\Person{撒提若}，在主教为其所作的墓志铭中，还加上了\Person{乌拉尼乌斯}一名。不过，这可能只是指他从尘世迁往天堂。\par

\Person{撒提若}早年与\Person{圣盎博罗削}一样，曾担任律师并任职。但当他的哥哥被任命为\Place{米兰}主教后，\Person{撒提若}立即放弃了自己的职务，献身于管理\Person{圣盎博罗削}的世俗事务，以免任何事分散他履行主教职责的精力。然而，在奉献于这项任务短短数年后，他于公元379年10月17日因重病去世。\par

\Person{圣盎博罗削}因失去这位志同道合的兄弟而悲痛欲绝，正因如此，我们才得以拥有在\Person{撒提若}葬礼上发表的精湛讲道，以及一周后论复活的第二次讲道。\par

\Person{圣盎博罗削}随后修订了这两篇讲道，它们以\Quote{圣盎博罗削论撒提若之死两卷}流传至今，一些手稿还补充道\Quote{与死人复活}。\par

据称由\Person{圣盎博罗削}为\Person{撒提若}所作的墓志铭如下：\newline{}Uranio Satyro supremum frater honorem\newline{}Martyris ad lævam detulit Ambrosius.\newline{}Hæc meriti merces, ut sacri sanguinis humor\newline{}Finitimas penetrans adluat exuvias.\par

% structure: p0008 source=h2 target=omitted reason=duplicate-page-title

1. 最亲爱的弟兄们，我们带来了我的祭献，一只无玷的祭献，一只悦乐天主的祭献，我的主和兄弟\Person{撒提若}。我没有忘记他是可朽的，我的情感也没有欺骗我，但恩宠却越发丰盈。因此，我无可抱怨，反而有理由感谢天主，因为我总希望，若有任何患难等待着教会或我自己，它们宁可落在我和我的家上。所以，感谢天主，在这普遍恐惧、因蛮族动向而万事堪忧的时刻，我以个人的悲伤结束了众人的烦恼，我为众人所恐惧的，都转向了我。愿此事完全成就，使我的悲伤成为众人悲伤的赎价。\par

2. 最亲爱的弟兄们，世间没有任何事物比这样一位兄弟更宝贵、更值得爱、更亲切，但公共事务优先于私人事务。若有人询问他的感受，他宁愿为他人被杀，也不愿为自己而活，因为基督按肉身为众人死了，好让我们学会不只为自己而活。\par

3. 此外，我无法对天主忘恩；因为我更应欣喜曾有这样一位兄弟，而非悲伤失去了兄弟，因为前者是恩赐，后者是当还的债务。因此，在我力所能及之时，我享受了委托给我的借贷，如今那寄存抵押品者已将其收回。否认曾收受抵押品，与因抵押品被交还而悲伤，并无区别。两者均显失信，而两者均危及永恒生命。拒绝偿还是一种过失，拒绝献祭则是虔诚。再者，金钱的出借人可能被愚弄，但自然的创造者、一切所需之物的出借人，却不能被欺骗。因此，借贷的金额越大，对资本的使用就越当感恩。\par

4. 所以，我无法对我的兄弟忘恩，因为他已归还原属于本性的东西，而获得了唯独属于恩宠的。谁肯拒绝共同的命运？谁会对特别委托给自己的抵押品被取走而悲伤？既然圣父已为我们交出了祂的独生子，任其死亡？谁认为自己应免于死亡的命运，既然他未能免于出生的命运？这是神圣之爱的一大奥秘：连在基督身上，肉身的死亡也未被豁免；祂虽是自然的主宰，却不拒绝自己所取肉身的律法。我必须死，而祂本不必。那曾对祂的仆人说：\Quote{我若要他等到我来，与你何干？}\BibleRef{若 21:22}的祂，难道不能若愿意就保持原样吗？但我兄弟在世上的延续，会毁掉他的赏报和我的祭献。还有什么比基督也按肉身死了更能安慰我们呢？或者我为何要为我的兄弟过度哀哭，既然我知道那神圣之爱不会死去？\par

5. 为何唯独我比众人更为他哀哭，而你们都在为他哀哭？我将个人悲伤融入众人的悲伤，尤其因为我的眼泪毫无用处，而你们的眼泪坚固信德，带来安慰。富人们哀哭，哭泣证明积聚的财富对安全无济于事，因为死亡无法用金钱推迟，最后一日同样带走富人和穷人。老年人哀哭，因为他们从\Person{撒提若}身上仿佛看见自己子女的命运；因此，既然不能延长肉身寿命，就当教导子女不追求肉体享乐，而追求德行职责。年轻人也哀哭，因为生命的终结不在乎年老。穷人也哀哭，而且更有价值、更富成果的是，他们以泪水洗刷了他的过犯。那是赎罪的泪水，那是隐藏死亡悲伤的叹息，这悲伤因永恒喜乐的丰盈而覆盖了先前悲伤的感觉。因此，虽是私人葬礼，哀悼却是公开的；所以，那被众人情感圣化的哭泣，不能长久持续。\par

因为，我最亲爱的兄弟，我为何要为你哭泣？你从我这被夺走，是为了成为众人的兄弟。我并未失去与你的交往，而是改变了方式；从前我们在肉体上形影不离，如今我们在情感上永不分离；因为你仍与我同在，并将永远与我同在。确实，当你与我同住时，故乡未曾将你从我夺走，你也从未将故乡置于我之前；如今你成了那另一个国土的担保，因为我在那里已不再陌生，我更好的那一部分已在那里。我从未完全沉浸于自身，我们各自的大部分都在对方内；然而我们每一个人都在基督内，在祂内有一切的全部，和每个人各自的份。这坟墓比你的出生地更令我喜悦，因为其中所蕴藏的并非本性的果实，而是恩宠的果实；在这如今无生命的身体里，安放着我生命更美好的成果，因为在我所负的这身体里，也有你更丰富的一部分。\par

但愿，正如记忆与感恩都归于你，我在这空气中尚存的一呼一吸，也能注入你的生命，并愿我的一半光阴能从我身上减去，加给你的！因为对于那始终共享祖产而不分彼此的人，生命之期本不应分割；或者至少，我们这些在世时总是毫无分别共享一切的人，在死亡上不应有所差别。\par

然而如今，兄弟，我该往何处去？我将转向何方？公牛寻找它的同伴，若失其伴，便自以为不全，以频频低鸣表达柔弱的渴念。而我，我的兄弟，难道不该思念你吗？我岂能忘记你？我曾与你同负这生命之犁。在劳作上我虽不及，但在爱中却更为紧密；并非因我的力量，而是因你的忍耐让我得以持久，你以焦虑的爱护之心，时常用你的身侧护着我的身侧，如兄弟般爱戴，如父亲般照料，如长者般警醒，如幼者般尊敬。你在同一亲等中为我尽了许多职责，以致我在你内失去的不仅是一个，而是许多；在你内，谄媚无从立足，孝爱历历在目。因为你无需以伪装增添什么，因一切已尽在你的孝爱之中，既无添加，亦无变更。\par

但我在过度的悲伤中走向何方？我忘记了本分，却记挂着所受的恩惠。宗徒将我叫回，仿佛给我的忧伤套上辔头，就如你们刚才所听到的：\Quote{我们不愿你们不知道，弟兄们，关于长眠的人，免得你们忧伤，像那些没有望德的人一样。}\BibleRef{得前 4:14} 请宽恕我，最亲爱的弟兄们。因为我们并非都能说：\Quote{你们当效法我，如同我效法基督一样。}\BibleRef{格前 4:16} 但若你们寻求可效法者，你们有一位可以效法。并非人人都适宜教导，但愿人人都善于学习。\par

但我们并未因流泪而犯什么大罪。并非所有哭泣都出自不信或软弱。自然的悲伤是一回事，不信的忧愁是另一回事；思念所失之物与哀叹失去之物的差别极大。不仅忧伤有泪，喜乐也有其泪。孝爱催人哭泣，祈祷浸湿床榻，恳求如先知所言，洗涤床铺。圣祖们安葬时，他们的朋友也曾大大哀哭。所以，泪水是热忱的标记，而非忧伤的起因。那么，我承认，我也曾哭泣，但主也曾哭泣。祂为一位非亲非故者哭泣，我为我的兄弟哭泣。祂在哭泣一人时哭泣了众人，我将为你而在众人中哭泣，我的兄弟。\par

祂为那触动我们而非祂自身的事哭泣；因为天主性不流泪；但祂在祂忧伤的那个本性内哭泣；在祂被钉死的那个本性内哭泣；在祂死去的、被埋葬的那个本性内哭泣。祂在今天先知所提醒我们的那个本性内哭泣：\Quote{熙雍母亲要说：一个人，是的，一个人在她内受造，至高者亲自建立了她。}祂在那本性内哭泣，在其中祂称熙雍为母亲，生于犹太，由童贞女怀孕。但按照祂的天主性，祂不可能有母亲，因为祂是祂母亲的创造者。就祂受造而言，并非由于神圣的出生，而是由于人的生育，因为祂成了人，天主降生了。\par

但你在另一处读到：\Quote{有一婴孩为我们诞生，有一儿子赐给了我们。}\BibleRef{依 9:6} 在\Quote{婴孩}一词中暗示年龄，在\Quote{儿子}一词中则含天主性的圆满。由祂的母亲所生，由圣父所生，但同一者既诞生又赐予，你不可当作两位，而是一位。因为天主子是独一的，既由圣父所生，又由童贞女所出，在次序上不同，但在名号上合而为一。正如刚才所读的训诲所教导的：\Quote{一个人在她内受造，至高者亲自建立了她；}的确，身体是个人，至高者是有权力。尽管祂在本质的差异上是天主又是人，但祂同时在两个本性内是一个。所以，一个特性专属祂自己的本性，另一个特性祂与我们共有，但在两者中祂是同一个，并且在两者中祂是完美的。\par

因此，天主使祂成为主和基督，这并不奇怪。祂使祂成为耶稣，即那在肉体本性中接受名字的祂；祂使祂成为圣祖达味所写的那一位：\Quote{熙雍母亲要说：一个人，是的，一个人在她内受造。}但既然成了人，祂与圣父不同，不是在天主性上，而是在祂的身体上；祂与圣父并非分离，而是在职分上有别，在权力上结合，但在受难奥迹中分离。\par

14. 这主题之处理需要更多论据，藉以证明圣父的权威、圣子的特性，以及整个天主圣三的合一；但今天我承担的是安慰之职，而非讨论之职，尽管在安慰时，通常借着投入讨论而将心灵从忧伤中引开。但我更愿缓解忧伤，而非改变情感，好使渴念得以缓和，而非被催眠。因为我不愿离我的兄弟太远，而被其他思绪引开，既然这篇讲辞是如同为陪伴他而承担，使我能以情感追随离去的他，并以心神拥抱我眼中所见的人。因我乐于将我双眼的全部凝视定在他身上，以慈爱的亲昵环绕他；而我的心神却麻木，仿佛他并未失去，我仍能看见他临在；我想他并未死亡，因我尚未察觉我对他的服侍有所欠缺——那些我曾以毕生与每一次呼吸所奉献的服侍。\par

15. 那么，对于如此恩情与如此辛劳，我能回报什么？我曾立你，我的兄弟，为我的继承人；你却使我成为继承人；我希望留你为生者，你却留下了我。我为了回报你的恩情，为报答你的恩惠，曾许下心愿；如今我失去了心愿，却未失去你的恩惠。作为我自己继承人的继承人，我该做什么？我比自己的生命活得更久，我该做什么？这光仍照耀着我，我却不再与它共享，我该做什么？我能向你回报什么感谢、什么善功？你从我这里得到的只有泪水。也许，你因确知你的赏报，并不渴求那些泪水——那是我仅剩的一切。因为即使你仍活着时，你也禁止我哭泣，并表明我们的忧伤比你的死亡更使你痛苦。泪水被命令不再流淌，哭泣被抑制。而对你感恩之情也禁止它们，免得我们在为损失哀哭时，似乎对你的功德绝望。\par

16. 但至少为我，你减轻了那忧伤之苦；我从前常为你恐惧，如今再无所惧。我没有任何东西是世界现在能从我手中夺走的。虽然我们圣洁的姊妹仍然健在，她一生无瑕，堪受尊敬，品性与你相等，在善行上也不逊色；然而我们俩从前更常为你担忧，我们觉得此生的一切甘甜都积聚在你身上。为你而活是乐事，为你而死不足悲，因为我们俩常祈求你幸存，而让我们活在你之后并非乐事。何时我们的灵魂本身不曾因这种恐惧的触动而战栗？你的疾病消息曾使我们的心神何等惊恐！\par

17. 哀哉我们可怜的希望！我们以为已得回所看见被夺去的人，如今才认出你的离去是因你向圣老楞佐殉道者所许的愿而获得的！我实在希望你不仅获得从此平安离世，也获得更长的寿命！你本可获得多年的寿命，既然你能获得从此离世。我实在感谢祢，全能永生的天主，祢至少没有拒绝给我们这最后的安慰，祢恩赐我们所爱的兄弟从西西里和阿非利加地区如愿归来；然而他归来后不久即被夺去，仿佛他的死亡只为这一点而被延迟，好使他能回到他的弟兄们身边。\par

18. 如今，我清楚拥有我的抵押品，任何变化都不能再将它夺去；我拥有可拥抱的遗骸，我拥有可用身体覆盖的坟墓，我拥有可躺卧其上的墓穴，且我将相信自己更蒙天主悦纳，因我将安息在那圣洁身体的骨骸上。但愿我曾能同样将自己的身体置于你的死亡之路上！若你曾遭刀剑攻击，我宁愿自己为你受刺；若我能在你生命逝去时唤回它，我宁愿献出自己的生命。\par

19. 接受你最后一息，或在你临终时向你口中呼出气息，对我并无益处，因为我以为要么我接受你的死亡，要么将我的生命转移给你。哦，最后一吻那可悲而甜蜜的抵押！哦，那拥抱的悲惨，其中无生命的身体开始僵硬，最后一息消散！我确实收紧双臂，但已失去所拥抱的他；我用嘴吸入你的最后一息，好分担你的死亡。但那气息在某种程度上成为给予我生命的，甚至在死亡中散发出更大爱的芬芳。若我未能以我的气息延长你的生命，至少愿你那最后一息的力量曾注入我的心灵，并使我们的亲情以你的纯洁与天真启发我。最亲爱的兄弟，你本会留给我这份遗产，它不会用悲伤的泪水打击情感，反而以卓越的恩宠举荐你的继承人。\par

20. 那么，如今我该做什么，既然那生命的一切甘甜、一切慰藉、最终一切魅力都已离我而去？因为你是我在家中的唯一慰藉，在外的唯一魅力；你，我说，是我的谋士、忧患的分享者、焦虑的驱除者、忧伤的赶逐者；你是我行为的保护者、思想的捍卫者；你，最后，是唯一肩负家庭与公共事务关怀的人。我呼求你圣洁的灵魂作证，在建造圣堂时，我常怕自己会使你不悦。最后，当你归来时，你责备自己的延迟。无论在家或外出，你都是司铎的指导者和教师，使他不容思考家事，而操心公共事务。但我不敢怕显得言过其实，因为这是你的赞美功绩：你，不惹任何人恼怒，既管理了兄弟的家庭，又举荐了他的司铎职。\par

21. 说实在的，我觉得我的心被你屡次的服务和你诸德的叙述所感动，虽然这样受感动，我却找到了安息；尽管这些记忆重新唤起我的悲伤，它们仍带来愉悦。难道我能不想你，或想到你而不流泪吗？难道我能够不记得这样一位兄弟，或者记得他而没有含泪的感激吗？因为什么事曾令我愉快，若不是源于你？我说，什么事曾使我喜欢而没有你，或使你愉快而没有我呢？我们不是几乎在每一件事上都共同实践，甚至我们的目光和睡眠也是如此吗？我们的意愿何时有过分歧？我们哪一步不是共同迈出的？以致我们几乎像是在抬脚时带动彼此的身体。\par

22. 但若有一方必须独自外出，人们都会认为他的一侧失去了保护，能看到他面容忧虑，会认为他的灵魂悲伤；往日的优雅、惯有的活力不再闪耀；这孤独对所有人来说都是一种可怕的事，令他们担心某种疾病。我们分离的事对所有人显得如此奇怪。我确实，因着兄弟的缺席而不耐烦，又时时想着此事，不住地转头寻找他，仿佛他就在眼前，而且那时似乎见到他并对他说话。但如果我希望落空，我就仿佛在自己低垂的颈上拖着一副轭，艰难前行，与别人相遇时满怀迟疑，匆匆回家，因为离开你再走远些就令我不悦。\par

23. 但当我们两人都必须外出时，路上的脚步并不比话语多，我们的步伐也不比谈话快；与其说是为了行走，不如说是为了交谈的愉悦，因为我们彼此都悬在对方的嘴唇上。我们并不留意凝视路过的风景，而是倾听彼此关切的谈话，饮入目光中善意的表情，吸入兄弟容颜的欢愉。我曾如何在内心中默默钦佩你的德行，如何庆幸天主赐给我这样一位兄弟——如此谦逊，如此能干，如此纯真，如此淳朴——以致当我想到你的纯真时，开始怀疑你的能力；当我看到你的能力时，几乎难以想象你的纯真！但你却以奇妙的完美将二者结合。\par

24. 最后，我们两人共同未能成就的事，你独自完成了。\Person{普罗珀尔}据我所知曾自庆幸，因为他以为因着我的司铎职分，他不必归还他所窃取的；但他发现你独自的力量大于我们两人合在一起的力量。于是他全部付清了，而且并未对你的节制忘恩负义，也没有嘲笑你的谦逊。但是，兄弟，你想为谁赢得这一切呢？我们希望那作为你劳苦之报酬的，正是它们的证明。你完成了所有事，而当完成一切归来时，你——你一个人，你本应被置于众人之上——却被从我们身边夺走了；仿佛你推迟死亡是为了这个目的：好让你完成爱的职务，然后摘取能力的棕榈枝。\par

25. 最亲爱的兄弟，这世界的荣耀多么少令我们喜悦，因为它们使我们彼此分离！我们接受它们，并非因为获得它们值得向往，而是为了不显得有卑劣的掩饰。或许它们之所以被赐予我们，是为使我们——因你将早逝而粉碎我们的欢乐——学会没有彼此而生活。\par

26. 我确实认出我心中的不祥预感，当我常常重读我所写的内容时。我曾尽力劝阻你，兄弟，不要亲自前往\Place{阿非利加}，而愿你差遣他人。我害怕让你踏上那次旅程，把你托付给海浪，一种比平常更大的恐惧笼罩了我的心灵；但你安排了行程，料理了事务，而且据我所知，又将自己托付给海浪，乘坐一艘老旧漏水的船。因为你急于求速，便将谨慎搁置一旁；渴望向我行善，你不顾自身的危险。\par

27. 啊，骗人的喜乐！啊，尘世之事无常的历程！我们以为那位已从\Place{阿非利加}回来、从海中得救、在船难后获保的人，如今不可能被从我们手中夺走；但是，虽然在陆地上，我们却遭遇了更悲惨的船难，因为那个大海的船难因着他强壮的游泳未能杀死的人，他的死亡对我们而言就是船难。因为还有什么快乐留给我们呢？从我们这里取走了如此可爱的装饰，如此明亮的光在这世界的黑暗中被熄灭了。因为在他身上，不仅是我们家族的装饰，而且是整个祖国的装饰都毁灭了。\par

28. 我确实对你们，最亲爱的弟兄们，圣洁的人们，怀着最深切的感激，因为你们把我的悲伤视为你们自己的悲伤，你们觉得这个丧失如同发生在你们自己身上，你们以非凡的爱予我全城的眼泪、各年龄段的哀悼以及各阶层的善意。因为这不是私人同情的悲伤，而是仿佛公共善意的服务与奉献。倘若因我失去如此一位兄弟而对你们的同情触动你们，我从中得到丰硕的果实，我得到你们爱心的保证。我宁愿我的兄弟还活着，但公共的善意在顺境中非常令人愉快，在逆境中也令人感激。\par

29. 于是，这样大的仁慈在我看来确实值得不平凡的感激。因为在《宗徒大事录》中，当\Person{塔彼达}死去时，寡妇们哭泣，\BibleRef{宗 9:39}或在《福音》中，群众因寡妇的眼泪而感动，并伴随那要复活的青年的葬礼。\BibleRef{路 7:12}那么，毫无疑问，由于你们的眼泪，宗徒的保护获得了；毫无疑问，我说，基督看见你们哭泣，就动了怜悯。虽然祂现在没有触摸棺架，但祂已接收了托付给祂的灵魂；如果祂没有用身体的声音呼唤死者，却以祂神性权能的权威，将我的兄弟的灵魂从死亡的痛苦和恶灵的攻击中解救出来。虽然死者没有在棺架上坐起来，但他已在基督内安息；如果他没有对我们说话，他却看见了那些在我们之上的事物，并因他现在看见比我们更高的事物而喜乐。因为从我们在福音中所读的，我们将理解将要发生的事，而我们目前所见是将来之事的预兆。\par

30. 他不需要为时间而复活，因为等待他的是为永恒的复活。因为为什么他该坠回这悲惨可怜的腐朽状态，回到这令人哀伤的生命？我们更应为他的灵魂从如此迫近的灾祸和威胁的危险中被拯救而欢喜。因为如果没有人哀悼\Person{哈诺客}，他在世界和平、战争未起时被提升，\BibleRef{创 5:24}人们反而为他庆幸，正如圣经论及他时说：\Quote{他被接去，免得恶念改变他的明悟，}\BibleRef{智 4:11}那么，如今当世上的危险加上生命的不确定时，更该怎样正义地说这话呢？他被接去，免得落入蛮族之手；他被接去，免得看见全地的毁灭、世界的终结、亲人的埋葬、同胞的死亡；最后，免得看见圣洁的童贞女和寡妇被玷污，这比任何死亡更痛苦。\par

31. 所以，兄弟，我认为你既因你生命的美丽，也因你死亡的适时而幸福。因为你被夺去，不是离开我们，而是离开危险；你没有失去生命，而是逃脱了威胁灾难的恐惧。因为你怀着圣洁心灵对身边之人的怜悯，如果你知道\Place{意大利}正因敌人的逼近而受压迫，你会怎样哀叹，怎样悲伤——我们的安全完全依赖\Place{阿尔卑斯山}的屏障，贞洁的保护在于树木的障碍！你会以怎样的痛苦哀悼你的朋友们被如此微小的分界线与敌人分开，而且是与一个既污秽又残忍的敌人，他不饶恕贞洁也不饶恕生命。\par

32. 我怎能说，你能忍受这些我们被迫忍受的事，并且或许（更痛苦的是）看见童贞女被强暴、幼童从父母怀抱中被夺走并抛在长矛上、献身于天主的身体被玷污，甚至年老的寡妇被污辱？我怎能说，你能忍受这些事？你甚至在临终时，忘记了自己，却仍不忘我们，警告我们关于蛮族的入侵，说你不是徒然说过我们应当逃命。或许是因为你看见我们因你的死亡而变得无助，你这样做，不是出于心灵的软弱，而是出于爱，你对我们是软弱的，但对自己是坚强的。因为当你被你的尊贵父母\Person{西马库斯}召唤回家，因为\Place{意大利}据说战火正燃，因为你正进入危险，因为你可能落入敌人手中时，你回答说，你来的原因正是为了在危险中不辜负我们，为了显示自己是你兄弟患难的同伴。\par

33. 他是幸福的，因着如此适时的死亡，因为他没有被保存下来经历这悲伤。当然你比你圣洁的姐妹更幸福，她失去了你的安慰，担忧自己的贞洁，最近还拥有两位兄弟，如今却因两者而痛苦，既不能跟随一个，也不能离开另一个；对于她，你的坟墓是旅店，你身体的安葬地是家。但愿这安息之所也是安全的！我们的食物伴着哭泣，我们的饮料伴着眼泪，因为你赐给我们眼泪的饼作为食物，并赐给我们大量的眼泪作为饮料，甚至过量。\par

34. 现在我自己该说什么呢？我不能死，以免离开我的姐妹，又不想活，以免与你分离。因为有什么能使我快乐，若没有你——你一直是我全部快乐所在？或者留在世上更久，在世上徘徊，还有什么满足？我们曾快乐地生活，只要我们在一起。如果这里有什么能使我们喜悦，没有你便不能喜悦；如果我们曾热切渴望延长生命，那么现在至少我们不会没有你而存在。\par

35. 这实在无法忍受。因为没有你——我生命中的同伴、我劳苦的分担者、我职务的共享者——还有什么可以忍受呢？我甚至不能通过预先思考来使他的损失变得更能忍受，因为我的心灵如此害怕想到任何关于他的这样的事！并不是我不知道他的状况，而是某种祈祷和誓愿蒙蔽了对共同软弱的感知，以至于我不知道如何思想关于他的事，除非是完满的幸福。\par

36. 然后最近，当我被一场严重的发作（但愿是致命的）所压倒时，我只悲伤你没有坐在我的床边，与我圣洁的姐妹分担友善的职责，在我死后用你的手指合上我的眼睛。我愿望了什么？我现在在思索什么？什么誓愿缺少？什么服务将要接续？我准备一事，却被迫陈述另一事；不是作为葬礼的主体，而是作为执行者。哦，坚硬的眼睛，能够看见我的兄弟死去！哦，残酷而不仁慈的手，合上了那双我从中看到如此多的眼睛！哦，更坚硬的颈项，能够承载如此悲伤的重担，尽管它是在充满安慰的服务中。\par

37. 我的兄弟，你原本更有理由为我做这些事。我本来指望从你手中得到这些服侍，我本来渴望它们。但现在，我比自己的生命活得更久，没有你，我能得到什么安慰？你曾是我悲伤时的安慰，是我快乐的激发者，是我忧愁的驱散者。我现在如何见到你呢，我的兄弟？你不再对我说话，不再给我亲吻？虽然，我们彼此的爱深深植根于我们心中，它更多是借着内在的情感而非公开的抚爱来滋养，因为我们既彼此信任和相爱，就不寻求他人的见证。我们兄弟情谊的强烈精神如此浸润了我们每个人，以至无需用抚爱来证明我们的爱；我们的心灵意识着彼此的情谊，满足于内在的爱，似乎不需要抚爱的表露，因为我们的外貌本身已塑造了彼此的爱；我们似乎，不知是由于何种精神的印刻或身体上的相似，彼此互为对方。\par

38. 谁看见了你，而不认为他已看见了我？我多少次问候那些人，他们因为先前已问候过你，就说已收到我的问候？多少人向你说了什么，然后述说他们已对我说过？这常常给我多大的快乐，多大的娱乐，因为我看见他们把我们弄错了？多么令人愉快的错误，多么可爱的疏失，多么无邪的欺骗，多么甜蜜的诡计！因为我在你的言语或行为中没有什么可恐惧的，而当它们被归于我时，我反倒欢喜。\par

39. 但如果他们过于强烈地坚持说曾告诉过我某些消息，我常常微笑着愉快地回答：\Quote{当心，你告知的不是我的兄弟。} 因为既然我们一切共有，同一精神，同一性情，然而唯独朋友的秘密并非公有财产，不是因为我们害怕沟通中的危险，而是为了因隐瞒而保持信实。但如果我们有需要商议的事，我们的商议总是共同的，尽管秘密并不总是公开的。因为虽然我们的朋友对我们任何一人说话，使得他们所说的可能传到另一人耳中；然而我知道，秘密多半被如此忠信地保守，以至连另一位兄弟也不被告知。因为这便是一个令人信服的证据：未被告知兄弟的事，也就没有被泄露。\par

40. 我承认，因着这些如此伟大卓越的恩惠而被提升到一种心灵的狂喜之中，我已不再害怕自己会成为幸存者，因为我认为他更配活着，因此我承受了这我无法忍受的打击，因为这种痛苦之伤在事先思想时比在意外来临时更容易承受。如今谁将安慰我这满腹忧伤的人？谁将扶起那被击倒的人？我将与谁分担我的忧虑？谁将使我摆脱这世俗的事务？因为你是我们事务的管理者，仆人们的监察者，兄弟与姊妹间的裁决者——不是在争执的事上，而是在情谊之事上裁决。\par

41. 因为如果任何时候我与我圣善的姊妹讨论某件事，哪个意见更可取，我们总是以你为仲裁者；你不会伤害任何人，又急于使每个人都满意，在裁决中保持你的友爱之情和公道，以便让每个人都满意地离去，并为自己赢得每个人的感谢。或者如果你自己提出某事来讨论，你的论证是多么愉快！而你的愤慨本身，又是多么不含苦涩！你的管教对仆人们本身来说也不是不愉快的！因为你宁愿在弟兄们面前责备自己，而非出于激动而惩罚！因为我们的修道身份抑制了我们纠正的热忱，而事实上，我的兄弟，当你承诺要惩罚却渴望减轻时，你消除了我们一切纠正的倾向。\par

42. 这乃是非凡审慎的明证，这种德性被智者如此定义：首要之善是认识天主，并以虔敬之心尊崇祂为真实而神圣的，以心灵的全部情感喜爱那永恒真理的可爱与可欲之美。其次在于从那神圣而属天的自然之源，生发出对近人的爱，因为连这世上的智者也从我们的律法借用了这些。因为若非从那神圣律法的天上源泉，他们绝不可能获得那些训诲人的要点。\par

43. 那么，关于他在崇拜天主一事上的虔敬，我说什么呢？在他被引入更完善的奥迹之前，遇到船难，他所乘的船撞上多石的浅滩，被抛来抛去的波涛撞碎；他不怕死，只怕未领受这奥迹就离开此世，于是向他所知道已领受入门圣事的人询问信者的神圣圣事；并非出于好奇想窥视奥秘之事，而是为获得信德的助佑。因为他让人将圣事用一块手巾包裹，将手巾系在颈上，这样投身入海，不寻求从船体上脱落的木板，借以浮在水上得救，因为他只寻求信德的工具。如此，他相信借此已得充分的保护与防御，便不再寻求其他帮助。\par

44. 同时也可以想一想他的勇气：当船碎裂时，他并非像一个遇难者那样抓住一块木板，而是像一个勇敢的人从自身找到了勇气的支撑，他的希望没有落空，他的期待也没有欺骗他。然后，当被波涛保存并安全地带到港口陆地时，他首先认出了他所托付的那位引领者——天主，立刻在亲自营救仆人或看到他们获救之后，不顾自己的财物，也不留恋失去的东西，去寻找天主的教会，好为他的获救而感恩，并承认永恒的奥迹，声称没有比感恩更大的义务。但如果对人不知感恩已被视为等同于杀人，那么对天主不知感恩是多么巨大的罪行！\par

45. 明哲之人，贵在自知，亦如智者所定，生活当合乎本性。然则，何事如此合乎本性，如同感恩于造物主？试观此天，当其被瞻仰时，岂非感谢其造主？因为\BibleQuote{高天陈述天主的光荣，穹苍宣扬祂手的化工。}海洋本身，当它平静安息时，便呈现神圣安宁的肖像；当它激荡时，则显示上界的义怒何等可畏。当我们看见无知无觉的大自然，仿佛有知有觉地抑制其波涛，波浪也知道自己的界限时，我们岂不都正当地钦佩天主的恩宠？我又该如何谈论大地？它服从天主的命令，慷慨地供给一切生物食物；田野则倍增其所接受的，犹如通过累积的利息而归还，并堆积起来。\par

46. 因此，他藉着本性的引导，以心灵的热切活力领会了神圣工程的法则，知道感恩首先应归于万有的保存者；但既然他无法回报，至少可以心怀感激。因为这感恩的本质在于：当它被献上时，它被感受到；而通过被感受，它得以献上。于是他献上感恩，并带回了信德。因为那位曾如此体验用巾帕包裹的天上奥迹的护佑的人，若他以口领受并引入内心深处，他会期待多少呢？当他将之接纳入怀时，对于那曾以巾帕覆盖即如此有益于他的奥迹，他岂不更期待更多吗？\par

47. 但他并非如此急切以至于放弃谨慎。他将主教召来，认为除非出于真信德，否则没有真正的感恩；他询问主教是否与天主教的诸位主教，即与罗马教会保持一致？或许在那地方，该地区的教会正处在分裂之中。因为那时\Person{卢齐斐尔}已退出我们的共融，虽然他因信德曾被放逐，并留下了他自己信德的继承者，然而我的兄弟认为，在分裂中不可能有真信德。因为尽管分裂者保持对天主的信德，但他们未保持对天主教会的信德，他们任由教会的某些肢体仿佛被分割，她的成员被撕裂。既然基督为教会受苦，而教会是基督的身体，那么那些使基督的苦难失效、分裂祂身体的人，似乎并未显示出对基督的信德。\par

48. 因此，尽管他保存了信德的宝藏，且因欠下如此巨债而惧怕航行，他仍宁愿渡海到一个可以安全偿还的地方，因为他深信对天主的感恩在于心志和信德；一旦他得以自由接近教会，他便毫不迟延地偿还了这感恩。他接受了所渴望的天主恩宠，并保全了所接受的。如此，没有什么比那能分辨神圣与人事的明智更为智慧的了。\par

49. 我又何必提及他在法庭职责中众所周知的辩才？他在最高长官的审判厅中引起了何等令人难以置信的赞叹！但我宁愿谈论那些他因默想天主的奥迹而认为比人事更可取的事物。\par

50. 若有人愿更充分地审视他的刚毅，请他想一想：他在遭遇海难后，如何以不可征服的轻生之态，多次渡海并旅行于广阔区域；最后，就在此时，他并未躲避危险，反而迎头而上。他忍受不义，不顾寒冷——但愿他同样谨慎防范！但正是因此他是有福的：只要体力允许，他度着实现青年事业的生活，不间断地完成他所愿做的，而不留意自己的软弱。\par

51. 然而，我怎能用言语表述他的纯朴？我指的是性情中的某种节制与心灵的清醒。请原谅，我恳求你们，将其归因于我的悲伤，若我允许自己较为充分地谈论那位我已不得与之交谈的人。当然，你们看到自己履行了这项友好的职责，并非出于软弱的感情，而是基于健全的判断；并非由于对他的死亡感到怜悯，而是出于尊崇他德行的愿望，这对你们是有益的；因为每个纯朴的灵魂都是蒙福的。他的纯朴如此伟大，仿佛变成了孩童，他以那无伪年龄的纯朴、完美德行的肖像、以及如镜子般反映的无辜品性而著称。因此他进入了天国，因为他相信天主的话，因为他像孩童一样拒绝了谄媚的诡计，而宁愿以温和承受不义的痛苦，也不愿严厉报复；他更愿意倾听抱怨而非诡诈，乐于和解，不慕虚荣，在谦逊上是圣洁的，以致在他身上，人更愿说有羞怯过度，而不必寻找所需的羞怯。\par

52. 但德行的基础从未过度，因为谦逊并不妨碍，反而有助于职责的履行。他的面容洋溢着某种童贞的羞怯，将内心的情感显露在脸上。若他突然遇见某位女亲属，他仿佛俯首并沉入地下；即使在男性同伴中也不例外，他很少抬头，很少举目，很少说话；当他做其中一件事时，总带着一种心中的腼腆谦逊，他身体的贞洁也与之一致。因为他圣洁地保全了圣洗圣事的恩赐，身体纯洁，内心更为纯洁；他不仅害怕身体上的羞耻，也害怕言谈中的不洁；他认为在言语的纯洁中保持谦逊与在身体的贞洁中保持谦逊同样重要。\par

最终，他如此热爱贞洁，以致从未寻求娶妻，尽管在他身上，这不仅是贞洁的渴望，也是他对我们所怀爱德的恩宠。然而，他以奇妙的方式隐藏了自己对婚姻的感受，并避免一切夸耀；他如此谨慎地隐藏其感受，以至于即便我们向他施压，他也显得更像是推迟婚期，而非回避婚姻。因此，这是唯一一点他不信任其兄弟和姐妹之处，并非出于任何犹豫不决，而纯粹是出于有德之谦逊。\par

那么，谁能不感到惊奇：一个人，在年龄上居于兄弟和姐妹之间（一个是贞女，一个是司铎），但在灵魂的伟大上不逊于二者，竟在两大恩赐上如此卓越，以至反映出一种圣召的贞洁和另一种圣召的圣洁，不是被职业所束缚，而是被德行的实践所束缚？如果情欲和愤怒生出其他恶习，那么我完全可以称贞洁和温和为德行之父母；尽管，正如它是一切善事的根源，虔诚也是其他德行的苗圃。\par

那么，我该说什么关于他的持家之道，一种对财产的节制呢？因为那些照料自己财物的人不寻求别人的财物，满足于自己所有的不因丰裕而骄傲。他除了自己的东西不愿追回任何东西，而且这样做与其说是为了更富有，不如说是为了防止被骗。他恰当地称那些寻求他人财物的人为\Quote{金钱之鹰}。但如果贪财是万恶之根\BibleRef{弟前 6:10}，不寻求金钱的人无疑已脱去了恶习。\par

他从不以精心准备的宴席或多道菜肴为乐，除非当他邀请朋友时，他愿望的是满足自然所需，而非为了享乐而过度丰盛。的确，他并非在财物上贫穷，而是在精神上贫穷\BibleRef{玛 5:3}。我们绝不应怀疑他的福乐，他既不作为富人因财富而欢喜，也不作为穷人觉得自己所有不足。\par

剩下的是，为达到诸枢德之终，我们应注意到他身上的正义之要素。因为尽管德行彼此关联且相连，但仍需要对每个德行作更清晰的描绘，尤其是正义。因为正义对自己有些吝啬，却完全致力于外在之事，凡它通过对自己的某种严格所拥有的，出于对众人的爱，它便倾注于邻人身上。\par

但这种德行有很多种：一种是面向朋友，另一种是面向众人，还有一种关于对天主的敬拜或对穷人的救济。那么他对众人如何，他所治理省份的人民之爱戴表明了；他们常说，他更像是他们的父母，而非判官；是慈爱诉讼者的仁慈仲裁者，是公正律法的坚定执行者。\par

但他在兄弟和姐妹面前如何，尽管他的善意包容了所有人，我们未分割的家产证明了，那份既未分配也未减少而是保存下来的遗产。因为他说爱不是立遗嘱的理由。他也在临终之言中表明了这一点，当时他称赞他所爱的人，说他的选择是永不娶妻，为不致与兄弟和姐妹分离；他也不立遗嘱，唯恐我们的感情在任何方面受伤害。最后，尽管我们恳切请求，他认为不应由自己决定任何事，然而他并未忘记穷人，只请求应给予他们在我们看来公正的数目。\par

仅凭这一点，他就充分证明了他对天主的敬畏，并树立了在人事方面虔诚的榜样。因为他给穷人的，是献给天主，因为\Quote{周济穷人的，是借给天主}\BibleRef{箴 19:17}；通过要求公正，他留给他们的不是一点点，而是全部。因为这就是正义的总体：变卖所有，分给穷人。因为那\Quote{散财，施舍给穷人，他的仁义永远常存}。所以他留给我们的是管家，而非继承人；因为对继承人而言，遗产是争议之事，而管家的职责是对穷人的义务。\par

因此，人可以正确地说，今天圣神藉着读经童子的声音告诉我们：\Quote{那手洁心清，不使自己的灵魂趋于虚妄，不对邻人使用欺诈的人，这是寻求上主的一代。}这样，他必要登上主的山，并居住在天主的帐幕中，因为\Quote{他行走无瑕，做了正直的事，说了真话，没有欺骗邻人}；他也没有把钱借给人生利钱，他总是希望仅仅保有继承下来的财产。\par

我何必再提他在虔诚上超越了单纯的正义：他认为考虑到我的职务，应给予我们财产的非法占有者一些东西，便宣布我是这番慷慨的发起人，但却将自己份额的收入转给了公共基金。\par

这些和其他事情，当时是令我心悦的，如今却加深了我悲伤的记忆。然而，它们固存，并将永远如此，绝不似影儿般消逝；因为德行的恩宠不随身体死亡，自然生命与功德也不同时结束，尽管自然生命的使用并非永远消亡，而是在一种暂时的豁免中安息。\par

那么，对于这样一个行了如此善事、脱离危险的人，我哭泣更多是由于对他的思念，而非因失去他。因为他死亡的及时性提醒我们，我们应以感激的敬意追随他，而非哀悼他；因为经上记载：个人的悲痛应在公共的哀悼中停止。这是用先知的语言说的，不仅指向那里所象征的那一位女子，也指向每个人，因为它似乎是向教会说的。\par

65. 那么，这话临到我，圣经说：\Quote{你教导这个，就这样训导天主的子民吗？难道你不知道你的榜样对他人是一种危险？除非你或许抱怨你的祈祷没有被垂听。首先，这是无耻的傲慢，渴望为自己获得你知道连许多圣人也被拒绝的事，而你清楚天主是不看情面的。} \BibleRef{宗 10:34} 因为虽然天主是仁慈的，但如果祂总是垂听所有人，祂就显得不再按自己的自由意志行事，而是出于某种必然。那么，既然所有人都祈求，如果祂垂听所有人，就没有人会死。因为你每天祈祷多少次？难道天主的安排要因你而失效吗？那么，你为什么哀叹有时得不到你知道不能总得到的事呢？\par

66. \Quote{愚昧的人，} 它说，\Quote{你在众妇女中，难道看不见我们的哀悼，以及发生在我们身上的事，即我们的母亲熙雍是如何充满忧伤而哀痛，如何因屈辱而自卑。现今你也要极其哀哭，因为我们都在哀哭；要忧愁，因为我们都在忧愁；而你为一个兄弟忧伤。去问大地，她必告诉你，她才是应当哀悼的，因为她活过了许多在她身上生长的人。从她而生，} 它说，\Quote{起初万物都从她而生，将来也要从她而出，看啊，他们几乎都走向灭亡，其中一大群被彻底拔除。那么，谁比失去了如此众多群羊的她更应当哀悼呢？而不是你，只为一个人伤心。}\par

67. 那么，让共同的哀悼吞没我们的哀悼，切断我们私人忧伤的苦涩。因为我们不该为那些我们看为得释放的人悲伤，并且我们记得，如此多的圣善灵魂此时此刻离开身体的锁链并非没有目的。因为我们可以看到，仿佛是出于天主的旨意，如此可敬的寡妇们接连去世，似乎是一种旅行的出发，而非死亡中的沉没，以免她们在侍奉天主满期中的贞洁陷入危险。多么悲痛的回忆在我心中激起怎样的呻吟和哀悼！而即使我没有空闲哀悼，但在自身个人的悲伤中，在如此多功绩的花朵的损失中，自然的共同命运安慰了我；我的悲伤因考虑一个人而遮掩了公开葬礼的苦涩，在家中显出虔敬的表象。\par

68. 那么，神圣的圣经啊，我再次寻求你的安慰，因为我喜欢默想你的诫命和你的语句。天地废去比法律的一笔一画落空更容易得多！但让我们现在听听所记载的话：\Quote{现今，} 它说，\Quote{你把你的悲伤留给自己，勇敢地承担所临到你的事。因为如果你承认天主的决定是公义的，你不但会按时接回你的儿子，而且在妇女中受赞扬。} 这话若是向一个女人说的，更何况向一位司铎呢！如果这样的话是关于一个儿子说的，那么用来谈论失去一个兄弟也并非不合适；尽管若他是我的儿子，我也无法爱他更多。因为正如在子女的死亡中，失去的劳苦和无益承受的痛苦似乎增加了悲伤；同样，在兄弟的情形中，交往的习惯和共同的职事加剧了悲伤的苦涩。\par

69. 但是，看！我听见圣经说：\Quote{不要再讲这话，却要听从劝告。因为熙雍的灾祸何其大！为耶路撒冷的忧伤得安慰吧。因为你看见我们的圣地被玷污，我们被呼求的名几乎被亵渎，我们的人遭受羞辱，我们的司铎被焚烧，我们的肋未人被掳去，我们的妻子被玷污，我们的处女受强暴，我们的义人被掳走，我们的幼童被交出，我们的青年被捆绑，我们的壮士变软弱。而其中最大的，是熙雍的印记失去了她的荣耀，如今她被交在我们仇敌的手中。那么，你要摆脱你的沉重忧伤，除掉众多的愁苦，好使大能者再次怜悯你，至高者藉减轻你的劳苦赐你安息。}\par

70. 于是，我的眼泪将要止住，因为人必须顺从有益的良药，因为信徒与不信者之间应当有所区别。让他们哭泣吧，那些没有复活希望的人——这希望并非因天主的判决，而是因信仰的严格而使他们失去。愿基督的仆人与偶像崇拜者之间有这个区别：后者为他们的朋友哭泣，认为他们永远灭亡；他们永不停止流泪，从忧伤中得不到安息，因为他们认为死者不得安息。但我们，对我们来说，死亡不是我们本性的终结，而只是今生的终结；因为我们本性本身被恢复到一个更好的状态，让死亡的来临擦去所有眼泪吧。\par

71. 而且，确实，如果那些认为死亡是感觉的终结和本性的消亡的人曾得到过任何安慰，那么我们更应当如此，因为对我们来说，善行的意识带来了更好赏报的应许！外邦人有他们的安慰，因为他们认为死亡是一切邪恶的止息；既然他们没有生命的果实，他们也认为他们逃脱了所有那些严酷、持续的痛苦的感受和折磨，即我们在今生必须承受的。然而，我们既然因赏报而更得支持，那么我们也应当因安慰而更忍耐；因为我们看他们不是失去，而是先被遣送，死亡不会吞噬他们，而是永恒要接纳他们。\par

72. 因此，我的眼泪应当停止，或者，若它们不能停止，我将为着你，我的兄弟，在共同的哀伤中哭泣，并将我私下的呻吟隐藏在公众的悲痛中。因为我的眼泪怎能完全停止呢？每当提及你的名字，或每当我的习惯动作唤起对你的记忆，或每当我的情感描绘你的形象，或每当回忆重新燃起我的悲伤时，它们便涌流而出。因为你怎么可能缺席呢？你在如此多的事务中重新呈现。我再说，你是现前的，总被带到我的面前，我以我的全心和灵魂拥抱你、凝视你、对你说话、亲吻你；无论是在幽暗的夜晚还是明亮的光照中，当你屈尊回访并安慰悲伤的我时，我捉住你。如今，那些在你生前曾令人厌烦的夜晚——因为它们剥夺了我们彼此相视的能力——以及睡眠本身，最近还曾是我们交谈的可憎打断者，已开始变得甜蜜，因为它们将你归还给我。那么，那些彼此同在永不中断、彼此关怀不减少、彼此尊重增长的人，不是痛苦的，而是有福的。因为睡眠是死亡的肖像和形象。\par

73. 但是，如果在夜的宁静中，我们的灵魂仍紧贴着身体的锁链，仿佛被囚禁在肢体的牢栏之内，却能够看见更高和更超脱的事物，那么，当它们处于纯净而属天的感官中，不受身体软弱的任何阻碍时，它们看见这些事物又该有多么多呢？因此，当某个黄昏临近时，我正在抱怨你没有在我安息时回访我，而你却始终完全临在。因此，当我四肢浸泡在睡眠中，我在心意上为你醒着时，你对我是活着的，我可以说：\Quote{我的兄弟，死亡是什么？} 因为毫无疑问，你从未与我分离片刻，因为你如此处处与我同在，以至于我们此生交往中无法享有的彼此享受，如今始终且处处与我们同在。因为在那个时候，当然并非一切事物都能临在，因为我们的身体构造不允许，也不能在任何时候、任何地方享受彼此相见、身体拥抱的甘饴。但我们灵魂中的图像始终与我们同在，即使我们不在一起时也是如此，这些图像并未终止，而是不断回到我们身边，渴望越大，我们所拥有的图像就越丰富。\par

74. 因此，我抓住你，我的兄弟，死亡和时间都不能将你从我身边夺走。眼泪本身是甘甜的，哭泣本身是一种愉悦，因为通过这些，灵魂的热切得以缓解，情感得以舒缓而平息。因为我不能没有你，也永不能忘记你，或想到你而不流泪。苦涩的日子啊！你显示我们的结合已被打破。值得流泪的夜晚啊！你为我失去了如此好的安息伴侣，如此难分难舍的同伴！你将使我遭受何等痛苦，除非那临在者的肖像呈现在我面前，除非我灵魂的异象代表他——我肉体的视觉不再显示的那一位。\par

75. 如今，如今，我的兄弟，我灵魂中最亲爱的，虽然你因过早的死亡而离去，至少你是有福的，因你不承受这些悲伤，也无需哀悼一位兄弟的失去——与他分离你本不能长久忍受，却迅速返回并再次探访他。但是，如果那时你急忙驱散我孤独的疲惫，减轻你兄弟心中的悲伤，那么如今你该多么常来探访我受折磨的灵魂，并亲自减轻那源于你的悲伤！\par

76. 但是，我职务的义务现在命令我稍事休息，对司铎职务的关注将我的心思引开；但我的圣善姐妹将会怎样呢？她虽然以对天主的敬畏节制自己的情感，却又以她虔敬的热忱再次点燃了情感本身的悲伤。她俯伏在地，拥抱她兄弟的坟墓，因劳苦行走而疲惫，心灵忧伤，日日夜夜更新她的悲伤。因为虽然她时常以言语中断哭泣，却在祈祷中重新哭泣；并且，虽然她在对圣经的认识上胜过那些带来安慰的人，她却以恒常的祈祷弥补她哭泣的渴望，主要在无人能打断她时，重新涌流大量的泪水。因此，你有可以怜悯之处，而非可以责备之处，因为在祈祷中哭泣是德行的标记。虽然这事在贞女中是常见的——她们较柔弱的性别和更温存的情感在看到共同的软弱时，即使没有家庭悲伤的感觉，也泪流满面——但当有更大的悲哀原因时，那悲哀就没有限度。\par

77. 安慰的方法，既然借口众多，便有所欠缺。因为你不能禁止你所教导的，尤其当她将眼泪归于虔敬而非悲伤，并因害羞隐藏共同悲伤的进程。因此，安慰她吧，你，你能够接近她的灵魂，渗透她的心思。让她察觉到你就在眼前，感到你没有离去，既已享受他的慰藉——她确信他的功绩——她可以学会不为他过度悲伤，因为他曾警告她，他不应被哀悼。\par

78. 但是，我为何要拖延你，兄弟？为何我要等待我的演讲死去，仿佛与你一同埋葬？虽然你无生命身体的视像、体型及其留于此处的余留俊美和形貌安慰了眼睛，我不再拖延，让我们前往坟墓。但首先，在众人面前，我道出最后的告别，向你们宣布平安，并献上最后的亲吻。你先我们而去，前往那共同并等待所有人的家园，而且如今我超乎众人地渴慕它。为你曾与之同居的那位准备共同的居所，正如在此我们曾拥有一切共有之物，同样在那里，也让我们知道没有分割的权利。\par

79. 不要，我求你，迟迟不愿接纳那渴望你的人，不期待那追赶你的人，不帮助那匆忙的人；若我看似迟延太久，请召唤我。因为我们从未长久分离，你总是习惯于归来。既然你不能再回来，我将到你那里去；我应当回报你的恩惠，轮到我前来了。在我们的生活境况中，从未有过太大差异；无论是健康还是疾病，都是两人共有的，以致若一人染病，另一人也病倒，当一人开始康复，另一人也随之痊愈。我们是如何失去了我们的权利？这一次，我们同样共有疾病，死亡如何不为我们共有呢？\par

80. 如今，我向祢，全能的天主，交托这无罪的灵魂，向祢献上我的祭献；请恩赐又仁慈地接纳一位兄弟的礼物，一位司铎的奉献。我预先献上我自己的初祭。我带着这抵押来到祢面前，这不是金钱的抵押，而是生命的抵押，不要让我长久拖欠如此大的债。这不是兄弟之爱的普通利息，也不是自然的寻常过程，却因如此大的德行而增加。我能承受，倘若我很快被迫偿还它。\par

\SourceRecord{Translated by H. de Romestin, E. de Romestin and H.T.F. Duckworth.；Nicene and Post-Nicene Fathers, Second Series；Vol. 10.；Edited by Philip Schaff and Henry Wace.；Buffalo, NY: Christian Literature Publishing Co.,；1896.；<http://www.newadvent.org/fathers/34031.htm>.}

% structure: p0001 source=h1 target=section reason=page-title

\section{论\Person{撒提鲁}之死（卷二）}

% structure: p0002 source=h2 target=omitted reason=duplicate-page-title

\Emph{论复活的信仰}.\par

在前一卷中，我稍许宣泄了心中的渴望，免得用过于猛烈的药膏敷在灼热的伤口上，反而加重而非减轻痛苦。同时，由于我常常对着我的弟兄说话，将他置于眼前，让天然的情感稍稍流露，也不是不合时宜的；因为泪水能稍予满足，哭泣能带来安抚，而打击则使之麻木。因为情感的外在表达，其性质柔和而温存，不爱任何过分、任何严厉、任何坚硬的东西；忍耐乃是在承受而非抗拒中得以证明。\par

因此，既然先前，那悲伤的景象足以使一个弟兄心思旁移，因为它完全占据了他；如今，既然在第七日——未来安息的象征——我们重又来到墓前，则颇有益处的是，将我们的思绪稍从我的弟兄转开，转向对众人的普遍劝诫，并专注于此事；这样既不以我们全部心思依恋我的弟兄，免得情感压倒我们，也不忘却如此敬爱与功绩，而完全离弃他；而且，倘若他的死又成为今日演讲的主题，只会加增我们强烈悲痛的苦楚。\par

因此，至爱的弟兄们，我们提议，以自然的普遍进程来安慰自己，不要认为等候众人的事是什么难忍之事。所以，我们认为死亡不应受到哀悼；首先，因为它是众人共有的，也是众人应得的；其次，因为它使我们摆脱今生的苦难；最后，因为当我们如同睡眠一样，从今世的劳苦中安息时，一种更活泼的精力便倾注于我们。有什么忧愁是复活的\OriginalTerm{恩宠}{grace}所不能安慰的呢？有什么悲伤，是相信死亡中无一物毁灭的信念所不能驱除的呢？不但如此，由于死亡加速，许多事物得以免于毁灭。所以，至爱的弟兄们，情形将是：在我们将普遍劝诫转向对我的弟兄的情感时，不会显得离他太远，如果藉着复活的希望和未来荣耀的美好，即使在我们的谈论中，他也能为我们而重新活过来。\par

那么，我们就从这一点开始，指出我们所爱之人的离世，我们不应为之哀悼。因为明知是预定为众人之事，却为之悲恸，仿佛是一桩特别的灾祸，还有比这更荒谬的吗？这等于是把心灵提拔到人的境遇之上，不接受共通的法则，拒绝天性的契合，在血肉的心智中妄自尊大，而不承认肉身自身的限度。不认识自己是什么，却假装成不是自己者，还有什么比这更荒谬呢？或者，明知某事将要发生，一旦发生却无法忍受，还有什么比这更缺乏先见之明呢？自然本身召唤我们回来，并以一种她特有的安慰，将我们从这类悲伤中引开。因为有什么如此深切的哀悼，或如此沉痛的忧伤，其中心灵竟然从不稍得缓解呢？人性有一种特点：尽管人可能处在忧伤的境遇中，但只要他们还是人，有时候也会稍稍转移思绪，离开悲痛。\par

确实，据说有些部落在婴儿诞生时哀哭，而在人死时举行庆典，这并非没有理由；因为他们认为，那些进入此生之海的人应受哀哭，而那些摆脱了今世波涛与风暴的人，则理所当然地应以欢乐相伴。我们也忘了亡者的生日，却以节庆的庄严重纪念他们逝世的那一天。\par

因此，依照自然，我们不应顺从过度的悲伤，免得我们似乎要么为自己索求一种超乎寻常的自然优势，要么拒绝那共有的命运。因为死亡对众人一视同仁，贫者无别，富者不免。所以，虽然因着一人的罪，死亡便临到了众人；\BibleRef{罗 5:12}这样，我们便不至于拒绝承认祂也是死亡之源，正如我们不拒绝承认祂是我们族类的创造者；并且，如何因一人而死属我们，复活也如何属我们；我们不应拒绝这苦难，为能获得那恩赐。因为，正如我们所读到的，基督\BibleQuote{来，是为拯救那丧失了的}\BibleRef{路 19:10}并\BibleQuote{成为死人并活人的主}\BibleRef{罗 14:9}。在亚当内，我跌倒了；在亚当内，我被逐出乐园；在亚当内，我死了；假如主不在亚当内找到我，怎能召我回来呢？正如我在他内是罪人，如今在基督内我却成了义人。那么，如果死亡是众人的债，我们必须能够承受这偿付。但这个话题须留待后文讨论。\par

现在我们旨在证明，死亡不应引起过重的忧伤，因为自然本身拒绝这样做。因此，据说在\Place{吕基亚人}中有一条法律，命令那些沉溺于悲伤的男子身穿女装，因为他们认为哀悼在男子身上显得柔软而女性化。而且，那些为了\OriginalTerm{信德}{faith}、为了宗教、为了国家、为了正义的判断、以及对德行的追求，本应挺身就死的人，若在他人身上遭遇此类事情时悲痛太过，这是不一致的；因为如果合适的原因要求，他们自己也会去寻求死亡。因为，当某事发生在他人身上时，你若以太少忍耐去哀悼，又怎能不厌恶这同样的事降临你自身呢？放下你的忧伤吧，如果你能做得到；如果不能，就留给你自己吧。\par

8. 那么，所有的悲伤都应藏在心里或加以压抑吗？为什么不让理智而非时间来减轻人的忧愁？智慧岂不更能平息那时间的流逝将抹去的伤痛？此外，在我看来，对我们所哀悼之逝者的记忆缺乏应有的情感，这表现在：我们宁愿遗忘他们，而不愿让安慰减轻悲伤；宁愿逃避对他们的回忆，而不愿怀着感恩记念他们；我们害怕想起那些其形象在我们心中本应成为喜乐的人；我们对于逝者被接纳一事，宁可猜疑而不抱希望，并且认为我们所爱的人更可能受罚，而非永生的继承者。\par

9. 但你可能会说：我们失去了我们曾经所爱的人。这岂非我们自身、大地与元素的共同命运，即我们无法永远保留那暂时托付给我们的东西？大地在犁下呻吟，被雨水鞭打，受暴风袭击，因寒冷捆绑，被太阳灼烤，为要结出年年的果实；当它以各色花朵装饰自己后，却被剥夺并抢走其装饰。它有多少掠夺者啊！然而它并不抱怨其果实的失去，它生养它们本就是为了失去它们，此后它也不拒绝结出它记得将会被取走的果实。\par

10. 天空本身也并非总是闪耀着闪烁星辰的光球，它们如同冠冕装饰着天空。天空并非总是因黎明的光芒而渐渐明亮，或因太阳的光线而泛红；而是在持续的交替中，那世界最悦目的景象随着夜晚的潮湿寒冷而变暗。有什么比光更令人感激？有什么比太阳更令人愉快？两者每天都走到尽头；然而我们并不因为这些从我们面前逝去而苦恼，因为我们期待它们回来。从这些事中，你被教导，对于属于你的人，你该表现出怎样的忍耐。如果天上的事物从你面前逝去而不引起悲伤，为什么人的逝去就要哀悼呢？\par

11. 那么，让悲伤保持忍耐吧，让逆境中有顺境所要求的那种节制。如果不合宜过分喜乐，难道过分悲伤就合宜吗？因为在悲伤或对死亡的恐惧中缺乏节制，是极大的恶事。这恶事曾驱使多少人走向绞索，曾将剑放在多少人手中，让他们以此证明自己的疯狂：既不能忍受死亡，却又寻求死亡；将逃避之事当作救药来接受。正因为他们无法忍受和承受那合乎他们本性的事，便陷入了与愿望相反的事，永远与那些他们渴望追随的人分离。但这并不常见，因为本性本身加以约束，尽管疯狂驱使人向前。\par

12. 但妇女们惯于公开哀号，仿佛害怕她们的痛苦不为人知。她们故意穿着污秽的衣服，仿佛悲伤之情就在其中；她们用污垢弄湿蓬乱的头发；最后，在许多地方常有这样的事：她们撕破衣服，将衣裳裂为两半，以裸露来出卖端庄，仿佛因失去了那原是端庄之回报的东西，便准备牺牲端庄。于是，淫荡的眼睛被挑动，对那些裸露的肢体生出欲念，若肢体不裸露，便不会引起欲念。但愿那肮脏的衣服遮盖的是内心而非形体。内心的放荡常隐藏在哀伤的衣服之下，而衣着的不雅粗野被用作遮掩放荡心灵秘密的遮盖物。\par

13. 那保持端庄、不放弃坚贞的妇女，便以足够的虔诚哀悼丈夫。对逝者应尽的最佳义务，就是让他们活在我们的记忆中，并继续存在于我们的爱中。那显明贞洁的妇女并未失去丈夫，也未因婚姻而守寡，只要她没有改换丈夫的姓氏。你协助共同继承人时，也并未失去继承人，反而用可朽事物的后继者换来了永恒事物的分享者。你有代表继承人的一位，当将欠继承人的给予穷人，好叫有一人存留下来，不仅存活到父亲或母亲的晚年，也存活到你自己的生命尽头。你留给继任者的就更多了，倘若他的那份不是用于现世事物的奢侈，而是用于购买将来的事物。\par

14. 然而我们渴望那些我们失去的人。有两件事特别令我们痛苦：一是对逝者的思念，这也是我自身的经历；另一则是我们以为他们被剥夺了生命的甘美，被夺去了辛劳的果实。因为有一种爱的温柔感动，突然点燃情感，以至于其效果更在于抚慰痛苦，而非阻抑痛苦；因为思念所失去的人似乎是一种本分，于是在美德的表象下，软弱反而增加了。\par

15. 然而，你为什么认为，那位因所爱的人去服兵役、或担任某种职务、或因经商而渡海，而送他往外地去的妇女，理应比你更忍耐呢？你并非因偶然的决定或贪财，而是因自然律而被留下。但你说，重获他的希望已被断绝。仿佛任何人的归来是确定的！常常，在危险恐惧强烈的地方，疑虑更使人疲惫；害怕某事发生，比承受已知已发生的事更为沉重。因为前者增加恐惧的程度，后者则期盼着悲伤的终结。\par

16. 然而主人有权随意将奴隶转移到任何他们认为合适的地方。难道天主没有这权利吗？我们未被允许期待他们回来，但被允许跟随那些已先行的人。当然，寻常生命的短促，似乎既没有剥夺先行者太多，也没有让存留者延迟太久。\par

但若人无法减轻自己的悲伤，那么因我们的思念而希冀整个事态的进程被扰乱，岂不显得不合宜吗？情人的思念确实更为强烈，然而他们因顾及必然之事而有所节制；虽然他们因被离弃而忧伤，却通常不会哀恸，反倒因被遗弃而羞愧自己爱得太仓促。因此，在遗憾中的忍耐便更加彰显出来。\par

但对于那些认为逝者被剥夺了生命的甜美的人，我又能说什么呢？在今生充斥的苦涩与痛苦之中，不可能存在真正的甜美，这些苦涩与痛苦或源于肉身本身的软弱，或来自外界所发生之事的不适。因为我们对于更幸福境况的愿望，总是焦虑不安、悬而未决；我们在不确定中摇摆不定，我们的希望将可疑之事当作确实，将不便之事当作满意，将终将衰败之事当作稳固；我们在意志上既无力量，在愿望上也无确据。但若有任何事违背我们的心愿发生，我们便以为失落了，且因患难中的痛苦而崩溃，远胜于因顺境中的享受而欢欣。那么，那些从烦扰中解脱的人，又失去了什么好处呢？\par

我不怀疑，健康对我们的益处大于疾病对我们的损害。财富带来的喜悦多于贫穷的烦恼，儿女亲情中的满足大于失去他们的悲伤，青春比老年的悲哀更令人愉悦。实现愿望常令人厌倦，渴望已久之事反成遗憾；以至于人们因获得了不害怕获得的东西而忧伤。然而，有什么祖国、什么享乐，能补偿流放和其他惩罚的痛苦呢？因为即使我们拥有这些，享乐也会因不愿享用或害怕失去而削弱。\par

但假设有人未受伤害，免于忧伤，能无间断地享受整个人生历程中的一切乐趣，那么，被束缚于如此肉身的桎梏中、受肢体狭窄界限限制的灵魂，又能获得什么安慰呢？如果我们的血肉之躯逃避牢狱，厌恶一切剥夺它漫游能力的事物；当它实际上似乎总要向外伸展，凭借其微弱的听或看的力量去探知超越自身之物时，我们的灵魂岂不更渴望从那\OriginalTerm{肉身的牢狱}{prison-house of the body}中逃脱吗？它自由运动如同空气，去往我们所不知之处，又来自我们所不知之处。\par

然而，我们知道灵魂在肉身死后仍然存留，并且一旦脱离肉身的栅栏，它便能以清晰的目光注视那些先前居于肉身时无法看见的事物。我们可以通过那些在睡梦中看见不在眼前乃至天上事物之人的事例来判断这一点（他们的心灵，当肉身仿佛被埋葬在睡眠中时，上升至更高的事物，并将它们传达给肉身）。所以，如果死亡将我们从今世的苦难中解脱出来，它当然不是恶事，因为它恢复了自由，排除了痛苦。\par

此时正是论证死亡并非恶事的恰当之处，因为死亡是脱离一切苦难和一切恶事的避难所，是安全的避风港，是安息的海湾。因为在这生命中，有什么逆境是我们未曾经历的呢？我们未曾遭受什么狂风暴雨呢？我们未曾被什么不适所折磨呢？谁的功绩能得幸免呢？\par

圣祖\Person{以色列}逃离本国，被迫离开父亲、亲族和家园，\BibleRef{创 28:5}，他为女儿的受辱，\BibleRef{创 34:2}，和儿子的死亡而哀恸，他忍受饥荒，死后又失去自己的坟墓，因为他恳求将自己的骨骸迁走，\BibleRef{创 49:29}，以免在死亡中仍不得安息。\par

圣\Person{若瑟}经历了弟兄们的仇恨，嫉妒他之人的诡诈，为奴的服役，商人的辖制，女主人的纵欲，她丈夫的不知情，以及监狱的痛苦。\par

圣\Person{达味}失去了两个儿子；一个是乱伦者，\BibleRef{撒下 13:29}，另一个是弑亲者，\BibleRef{撒下 18:14}。拥有他们是耻辱，失去他们是悲痛。他又失去了第三个，是他所爱的婴孩。婴孩活着时，他为他哭泣，但死后却不再渴想。因为我们读到，当孩子患病时，达味为他恳求主，禁食并身穿苦衣躺卧；当长老们前来要从地上扶起他时，他既不愿起来也不进食。但当他听说孩子死了，就更换衣裳，敬拜天主，然后进食。当这事在他的仆人看来奇怪时，他回答说，孩子活着时他禁食哭泣是对的，因为他合理地认为天主或许会施怜悯，而那能赐生命给逝者的，无疑也能保全活者的生命；但如今死亡既已发生，他为何还要禁食呢？因为他如今已不能带回那死去的人，召回那已无生命的人。他说：\BibleQuote{我要到他那里去，他却不能回到我这里来。}\par

啊，为哀恸者带来多大的安慰！啊，智者的真判断！啊，身为奴仆者的奇妙智慧！没有人应当因遭遇逆境而恼怒，或抱怨自己受到与功劳不符的磨难。因为你是谁，竟预先声称自己的功劳？你为何想要抢先于洞悉万事的那一位？你为何要从将要审判的那一位手中夺去判决？就连圣人也不容许如此行，圣人也从未不受惩罚地如此行过。\Person{达味}在自己的圣咏中承认，他因此事而受鞭笞：\BibleQuote{看，这些不虔敬的人，在世享平安，积攒了财富。所以我徒然洗净了我的心，又在无辜者中洗了我的手；我终日受鞭笞，每天早晨都受责罚。}\par

27. 伯多禄也是如此，虽然充满信德和热忱，却因为尚未意识到我们共有的软弱，冒失地对主说：\BibleQuote{我要为你舍命}\BibleRef{若 13:37}，便在鸡叫三次之前陷入了冒失的试探\BibleRef{路 22:60}。虽然那试探确实是为我们的救恩而设的一课，好叫我们学会不可轻视肉身的软弱，免得因轻视它而受试探。如果伯多禄也受了试探，谁还能自以为站得稳呢？谁还能说自己不会受试探呢？无疑，伯多禄是为我们而受试探的，这试探的验证并不发生在比他更强的人身上，却要在他身上使我们学习如何在试探中抵抗，纵然受制于对生命的眷顾，也能以忍耐的泪水克胜试探的刺痛。\par

28. 不过，还是那位达味，为了不使那些固执圣经字句的人因他行为的差异而困惑，我再说，还是那位达味，未曾为无辜的婴孩哭泣过，却为死去的弑亲者哭泣了。因为最后，当他哀哭悲恸时，他说：\Quote{我儿阿贝沙隆！我儿阿贝沙隆！谁许我替你死了呢？}但哭泣的对象不单是弑亲者阿贝沙隆，还有阿默农；不单为这乱伦者哭泣，甚至还要为他报仇；一个因遭受王国的轻蔑，一个因兄弟被流放。受哭泣的是恶人，而非无辜者。原因何在？理由为何？谨慎之人的深思和智慧之人的印证不容轻视；因为行为如此殊异，却蕴含极大的明智一致，但信仰却是一体的。他为那些已死的人哭泣，却认为不应为那死去的婴孩哭泣，因为他觉得前两者对他而言是失丧了，却希望后者能复活。\par

29. 关于复活，后文将详加讨论；现在让我们回到当前的主题。我们已经阐明，即使圣人们，也未曾顾及自己的功德，在此世承受了许多沉重的苦难，伴随着劳苦与不幸。因此，达味自省时说：\BibleQuote{主，求你记念，我们原是尘土；世人的日子有如青草。}在另一处又说：\BibleQuote{人如同虚幻，他的日子消逝有如阴影。}还有谁比我们更可怜呢？我们被遣入此世，仿佛被剥夺净尽、赤身露体，带着脆弱的身体、狡诈的心、软弱的意志，为挂虑而焦虑，在劳作上怠惰，在享乐上放纵。\par

30. 按撒罗满的看法，未出生乃是最佳。就连那些自认为最精通哲学的人也追随了他的见解。他虽比这些哲学家时代更早，却比我们的许多作者为晚，在\WorkTitle{训道篇}中这样说：\BibleQuote{我称扬那些早已死去的死人，胜过那些还活着的活人。那还没有出生，没有见过太阳下所行的恶事的人，比这两种人更为有福。我又看见了一切劳苦，和一切工作的成就，都只是出于人彼此间的嫉妒：这也是虚幻，也是追风。}\par

31. 说这话的，除了那求得并获得了智慧，为知道世界如何造成，元素的能力，年岁的运行，星辰的排列，了解生物的本性，野兽的凶猛，风的强力，并洞察人的思念的人之外，还有谁呢？那么，那些连天上的事都不能隐藏于他的人，人世间的事又怎能隐藏于他呢？他洞悉那争夺他人孩子的妇人的心思，凭着圣宠的启示明了那些他未曾经历过的生物的本性；关于他自身经验所发现的、属于那本性的境况，他岂会犯错或说不实之言呢？\par

32. 但并非只有撒罗满一人有此感受，尽管只有他说了出来。他读过圣约伯的话：\BibleQuote{愿我出生的那天消灭}\BibleRef{约 3:3}。约伯认清了出生是一切灾祸的开端，因此他愿自己出生的那天消灭，好使一切烦恼的根源得以除去；他愿自己出生的那天消灭，为能领受复活的日子。因为撒罗满听过他父亲的话：\BibleQuote{主，求你使我认识我的终期，和我的寿数，好叫我明了我的缺欠如何。}因为达味知道，在此我们无法领悟圆满的，因此他急切奔赴那将要来到的事。因为我们如今是凭着镜子观看，模糊不清，到那时，就能面对面地看清了；如今我所知有限，到那时就全知道，如同主知道我一样。当那并非阴影而是天主尊威与永恒的真体开始照耀，且被我们以没有帕子遮蔽的脸凝视时，圆满才能被领悟\BibleRef{格前 13:12}。\par

33. 然而，除非是逃离此生的困苦，没有人会急于奔赴终点。因此达味也解释说，他为何急于奔赴终点：\BibleQuote{看，你使我的岁月屈指可数，我的生命在你面前如同无有；凡是活着的人，尽属虚幻。}那么，我们为何还迟疑逃离虚幻呢？为何还喜欢在这世上徒然烦扰，积蓄财宝，却不知道为哪个继承人积存呢？让我们祈求远离烦扰，使我们被接出这愚妄的世界，得以脱离每日的漂泊，回归那乡邦，我们本然的家乡。因为我们在这世上原是旅客和外方人；我们必须返回那来的地方，必须努力并恳切地祈祷，而非敷衍了事，以求脱离那些巧言令色者的诡诈和邪恶。那知晓补救之道的人，叹息自己寄居的时日延长，不得不与不义的人和罪人同居。我既是有罪的，又不知道补救之道，我又当如何呢？\par

耶肋米亚也曾以这些话哀叹自己的诞生说：\BibleQuote{我真不幸，我的母亲！你为何生了我这个在普世上与人争辩的人？我没有施惠于人，也没有人施惠于我；我的力量已经衰竭。} 如果圣善的人尚且畏避生命，他们的生命虽有益于我们，他们却视为于己无益；那么我们这些既不能利及他人，又感到自己的生命有如借来的钱子，每日不断加重的罪债使它日益沉重的人，又当如何呢？\par

\BibleQuote{我天天死。}\BibleRef{格前 15:31}宗徒如此说。这句话显然比那些声称默想死亡乃是真智慧的人所说的更好，因为他们称赞这种研究，而他却实践了死亡。他们单为自己而行，但保禄自己虽是成全的，他死却不是为自己的软弱，而是为我们的软弱。然而，默想死亡不是别的，正是灵魂与肉身的一种分离，因为死亡本身无非被定义为灵魂与肉身的分离吗？但这只是依常人之见。\par

然而，按照圣经，我们受教得知死亡有三种。一种死亡是当我们死于罪恶，却向天主活着。那真是有福的死亡，它摆脱罪恶，归向天主，使我们脱离可朽坏的，并祝圣我们归于那不死不灭者。另一种死亡是离开此生，如圣祖亚巴郎的死，和圣祖达味的死，他们与自己的祖先同葬；那时灵魂从肉身的束缚中被释放。第三种死亡便是那所说的：\BibleQuote{任凭死人去埋葬他们的死人罢！}\BibleRef{玛 8:22}在这死亡里，不独肉身死亡，连灵魂也死，因为\BibleQuote{犯罪的人，必死无疑。}\BibleRef{则 18:4}因为它向主而死，并非由于本性的软弱，而是由于罪过。但这种死亡并非脱离此生，而是因错谬而堕落。\par

这样，灵性的死亡是一回事，本性的死亡是另一回事，惩罚的死亡是第三回事。但那本性的死亡并非也是惩罚性的，因为主并没有将死亡作为惩罚施加于人，而是作为治疗。亚当犯罪后，有一件事被定为惩罚，另一件事则作为治疗，经上说：\BibleQuote{因为你听从了你妻子的声音，吃了我曾吩咐你唯独不可吃的那树上的果子；为了你的缘故，地成了可咒骂的；你一生日日劳苦，才能吃它的出产。地要给你生出荆棘和蒺藜，你要吃田间的蔬菜；你必须汗流满面，才有饭吃，直到你归于土中，因为你原是由土来的。}\par

你看，惩罚中的休息之日就在于此，因为它们包含着对今生的荆棘、世俗的挂虑，以及拒斥圣言的财富享乐所判处的刑罚。死亡被赐予作为治疗，因为它是诸恶的终结。因为天主并没有说：\BibleQuote{因为你听从了你妻子的声音，你必归于土中}——那会是惩罚性的判词，正如下面这句：\BibleQuote{受咒骂的地必给你生出荆棘和蒺藜}——而是说：\BibleQuote{你必须汗流满面，才有饭吃，直到你归于土中。}你看出死亡更是我们种种惩罚的终点，它给此生的行程画上句号。\par

所以，死亡非但不是恶，反而是某种善。因此它被视为善事而受人寻求，正如经上所载：\BibleQuote{人必求死，却不得死。}\BibleRef{默 9:6}那些要对山岭说：\BibleQuote{倒在我们身上罢！对丘陵说：遮盖我们罢！}\BibleRef{路 23:30}的人，也必寻求它。那犯了罪的灵魂，也必寻求它。那躺在地狱里的富人要寻求它，他想要拉匝禄用指尖凉润他的舌头。\BibleRef{路 16:24}\par

这样，我们看出，此死亡乃是利益，而生却是惩罚，所以保禄说：\BibleQuote{因为在我看来，生活原是基督，死亡乃是利益。}\BibleRef{斐 1:21}基督若不是身体的死亡，生命的气息，又是什么呢？因此，让我们与祂同死，好使我们能与祂同活。那么，在我们内心应当有一种每日的操练和趋向死亡的倾向，藉着这种与我们说过的肉身欲望的分离，我们的灵魂可以学会退隐，并如同被置于高处，当尘世的贪欲不再能靠近并依附灵魂时，灵魂便能承受死亡的模样，以免招致死亡的惩罚。因为肉身的法律与理智的法律交战，并将它转交于错谬的法律之下，正如宗徒晓示我们的，他说：\BibleQuote{可是我发觉，在我的肢体内，另有一条法律，与我理智所赞同的法律交战，并把我掳去，叫我隶属于那在我肢体内的罪恶的法律。}\BibleRef{罗 7:23}我们都被牵连，我们都感受到这点；但我们并非都得蒙解救。因此，我真是个不幸的人，除非我寻求那治疗。\par

但治疗是什么呢？\BibleQuote{谁能救我脱离这该死的肉身呢？感谢天主，藉着我们的主耶稣基督。}\BibleRef{罗 7:24}我们有一位医生，让我们使用这治疗。我们的治疗是基督的恩宠，而该死的肉身正是我们的身体。因此，让我们与自己的身体成为陌路人，以免我们与基督成为陌路人。我们虽在肉身内，却不要随从肉身的事，也不要弃绝本性的正当需求，而要在一切之上渴慕恩宠的礼物：\BibleQuote{我正夹在两者之间：我渴望求解脱而与基督同在一起，这实在是再好没有了；但存留在肉身内，对你们却更为重要。}\BibleRef{斐 1:23}\par

但主耶稣，这种需要并非人人的情形；对我而言并非如此，我对谁也不受益；因为我，死亡乃是利益，好使我不再犯罪。死亡为我乃是利益，我正撰写这篇论著安慰他人之际，自己却仿佛受到一股强烈冲动的催迫，切望我已失去的兄弟，因为这冲动不容我忘怀他。如今我爱他更深，也更强烈地切望他。我说话时切望他，重读写下的文字时也切望他，我想我愈受催迫写下这些，正是为使我永不至于没有对他的怀念。而我这样做，并不违反圣经，反而与圣经同心合意，好能以更大的忍耐哀恸，以更大的切望渴念。\par

43. 吾弟，你使我不再惧怕死亡；我只愿我的生命与你的生命同逝！这正是\Person{巴郎}因先知之神所感，为自己所祈求的最大福祉，他说道：\BibleQuote{愿我的灵魂死于义人的灵魂，愿我的后裔如同他们的后裔。}他确实是按先知之神如此祈求的，因为他既看见了基督的兴起，也看见了祂的胜利，看见了祂的死亡，同时也看见了在祂内人类永恒的复活，因此他不惧怕死亡，因为他将复活。那么，不要让我的灵魂死于罪恶，不要让罪恶进入它内，而要让它在义人的灵魂中死去，以承受他的正义。这样，在基督内死去的人，也在\OriginalTerm{圣洗}{Font}中分享祂的恩宠。\par

44. 因此，死亡既非可怖之物，对困乏者不苦涩，对富者不过于痛苦，对老者不残酷，对勇者非懦弱之标志，对忠信者非永恒，对智者亦非意外。因为多少人仅凭其死亡的声望就圣化了生命，多少人因羞于存活而以死为利！我们读到过，许多次，伟大民族的得救是借着一人的死亡；敌人军队因统帅之死而溃逃，而此统帅生前却未能征服他们。\par

45. 借着殉道者的死亡，宗教得以捍卫，信仰得以增长，教会得以坚固；死者得胜了，迫害者被克胜了。因此，我们庆祝那些我们对其生平一无所知者的死亡。同样，\Person{达味}也曾先知性地因自己灵魂的离去而欢欣，说：\BibleQuote{在上主的眼中，圣者的死是珍贵的。}他视死亡胜于生命。殉道者们的死亡本身正是他们生命的奖赏。再者，借着争执者的死亡，仇恨也得以终结。\par

46. 何必多言？世界借着一人的死亡而得救赎。因为基督，如果祂愿意，本无需死亡，但祂既不认为死亡应被逃避，仿佛其中有何懦弱，也不能以比死亡更好的方式拯救我们。因此，祂的死亡就是所有人的生命。我们以祂死亡的印记被标记，我们祈祷时宣示祂的死亡；我们献上\OriginalTerm{祭献}{Sacrifice}时宣告祂的死亡，因为祂的死亡是胜利，祂的死亡是我们的奥迹，祂的死亡是世界一年一度的隆重庆典。如今，关于祂的死亡，我们还能说什么呢？因为我们借这神圣的表样证明，唯有死亡获得了不死不灭，死亡自身救赎了自身。因此，死亡不应被哀悼，因为它是众人得救的缘由；死亡不应被逃避，因为天主子不以为耻，也未逃避它。本性的秩序不应被破坏，因为万民共通之事不能容许个人例外。\par

47. 事实上，死亡本不属于人的本性，却成了本性之事；因为天主起初并未设立死亡，而是将之作为补救之方赐下。因此我们当留意，别让它显得相反。因为若死亡是好的，为何经上记载说：\BibleQuote{天主并未造死亡，但因人的恶意，死亡进入了世界}呢？因为实际上，死亡并非天主化工的必然部分，因为对那些安置于乐园中的人，一切美物的连续洪流曾涌出；但由于过犯，人的生命被判要经受漫长的劳苦，开始带着不堪忍受的痛苦呻吟而困苦；因此适宜为诸恶设定一个终结，让死亡恢复生命所失落的。因为不死不灭，若非恩宠灌注，与其说是益处，倒不如说是负担。\par

48. 若仔细思考，这不是我们存有的死亡，而是邪恶的死亡，因为存有得以延续，邪恶却灭尽了。那已逝去的将复活；但愿它既如今已无罪愆，亦无有先前的罪过！但正是此一事实证明这并非存有的死亡，即我们将仍是我们从前之人。因而，我们将或是清偿我们罪罚，或是获得我们善行的赏报。因为同一存有将复活，且因偿付了死亡的代价而更为尊贵。然后\BibleQuote{在基督内死去的人，要先复活；}经上说，\BibleQuote{然后我们这些活着还存留的人，同时与他们一同被提到云彩之上，到空中迎接主：这样，我们就时常同主在一起。}\BibleRef{得前 4:16}他们先复活，而我们这些活着的人随后。他们偕同耶稣，那些活着的人借着耶稣。对他们而言，安息之后的生命将更为甘饴，尽管活着的人将享有喜乐的增益，他们却未曾体验过补救之方。\par

49. 因此，在死亡中我们无所惧怕，无所哀悼，无论从本性所领受的生命被归还给本性，还是为某种要求它的责任而牺牲，这都将成为宗教行为或某种德行的实践。从未有人愿意保持现状。人们曾认为这是许诺给了\Person{若望}，但事实并非如此。我们持守经文，并从中推导出意义。他在自己的著作中否认曾有他不死的许下，\BibleRef{若 21:23}以免有人因此事例而对虚妄的希望有所凭恃。但若连愿望此事都是过度的希望，那么对循例发生之事而无节制地哀伤，又该是多么过度呢！\par

50. 异教徒大多用以下思想来安慰自己：或是共同的悲苦，或是本性的法则，或是灵魂的不死不灭。但愿他们的言论能前后一致，不至于将可怜的灵魂变成许多可笑的怪物和形态！但我们这些以复活为赏报的人，又当如何呢？尽管许多人因无法否认这恩赐的伟大而拒绝相信它。为此我们将坚持此一信仰，不仅靠一偶然的论据，而是尽我们所能地据理力争。\par

51. 其实，一切事物之所以为人所信，或是出于经验，或是基于理智，或是通过类比，或是因其适宜存在，而每一种都支持我们的信仰。经验教导我们，我们是被推动的；理智则指出，推动我们的那一位必定属于另一种能力；类比则表明田地曾出产庄稼，因此我们期望它将继续出产。适宜性则在于，即使我们认为不会有结果，却依然相信放弃德行之工是绝不相宜的。\par

52. 那么，每一条理据都是相互支持的。然而，对复活的信仰最清楚地依据三点推断出来，而这三点已涵盖一切。这便是理智，从普遍现象中得出的类比，以及已发生之事的证据，因为已有许多人复活了。理智是清晰的。因为我们生命的整个过程在于灵魂与身体的结合，而复活所带来的，或是善工的报酬，或是恶行的惩罚，所以那行为被衡量的身体，必须复活。因为若要为灵魂与身体的共融交账，没有身体，灵魂如何能被传唤受审呢？\par

53. 复活是所有人的命运，但信此有困难，因为它并非源于我们的功劳，而是天主的恩赐。复活的第一论证，便是世界的运行，万物的状况，世代的接续，更迭中的变迁，星辰的起落，昼夜的终始，及其每日的交替，仿佛重新得生而出现。这大地之所以有肥沃的特性，别无其他原因，只因天主所定的秩序，借着夜间露水，将被白日太阳的炎热晒干的水分，重新补足，而地上万物皆由此水分而生。我何须再提地上的果实？它们落下时，岂不看似死去，而再度变绿时，又复生起来？所播种的，复活起来；已死亡的，又复活起来，并且它们重又形成与先前同类的物种。大地首先归还了这些果实，在这些果实中，我们的人性首先找到了复活的模式。\par

54. 何必怀疑身体将从身体复活呢？谷粒被播种，谷粒再长起；果实被播种，果实再长起；然而谷粒却披上花和壳。\BibleQuote{这可朽坏的，必须穿上不可朽坏的；这可死的，必须穿上不可死的。}\BibleRef{格前 15:53}复活的华美便是不朽，复活的华美便是不可朽坏。因为有什么比永远的安息更为丰盛？有什么比永久的平安更为富足？这便是那丰盛的果实，因着它的增长，人性在死后更为丰沛地萌发。\par

55. 然而你惊奇，那已经腐烂的东西怎能再度结实，分散的微粒怎能重聚，被消灭的怎能得以复原：你却未惊奇于种子如何在大地湿润的压力下破碎后犹能发绿。因为它们确实也在与土壤接触下腐烂、破碎，而当土壤肥沃的湿润赋予那些死去且隐藏的种子以生命，且借生命的热力，仿佛为草本植物呼出一种灵魂。随后，大自然一点点地从地里托起那生长谷穗的嫩茎，犹如一位细心的母亲，用某些叶鞘将它包裹，以免它在生长时被严寒的冰霜所伤，又保护它免于太阳的过分炎热；而此后，又怕雨水会打落那仿佛刚从最初摇篮中挣脱、刚刚长大的果实，或是风将它吹散，或是小鸟将它啄食，于是她通常用芒刺的篱笆将它围护起来。\par

56. 那么，大地既赋予它所接受的任何种子以生命，使之生长、穿戴、保护并防卫，如果它将所领受的人的身体也归还出来，何必惊奇呢？因此，不要再怀疑：那忠信的大地，既好像以利息般加倍归还所托付的种子，也必将归还所托付的人类之抵押品。我又何必再提及那各类树木呢？它们由所播之种萌发，以复苏的繁殖力再次结出初绽的果实，重复旧日的形状与形象，有些树木得以更新，延续多代，甚至穿越若干世纪而长存。我们见到葡萄腐烂，葡萄树却又长出来：嫁接插枝，树木便得再生。难道有这恢复树木的天主的眷顾，却对人没有关怀吗？那位不曾容许祂为人的享用而赐下的东西灭亡的，岂会容许祂按自己肖像所造的人灭亡吗？\par

57. 但你觉得死人复活是难以置信的吗？\BibleQuote{糊涂人哪！你所播的种子若不先死了，决不能生。}\BibleRef{格前 15:36}你随意播下任何干种子，它都被举扬起来。但你回答说，种子自身含有生命液。我们的身体也有血，也有自己的湿润。这便是我们身体的生命液。因此，我认为有人提出的反对意见就不攻自破了，他们说干枯的树枝不能再生，然后竭力借此来攻击肉身。因为肉身并非干枯，因为一切肉身属于泥土，泥土出于湿润——湿润来自大地。再者，许多生长的植物，虽常青翠，却出于干燥多沙的土壤，因为大地本身为自身提供足够的湿润。那么，这不断使万物复原的大地，岂会在人身上失效呢？由上述可见，我们不应怀疑这更合乎本性，而非违背本性；因为一切有生命之物复活是合乎本性的，而灭亡才是违背本性的。\par

我们现在来谈一个让外邦人深感困扰的问题：大地如何能恢复那些被大海吞没、被野兽撕裂或吞噬的人呢？于是，最终我们不得不做出结论：怀疑并非针对普遍的复活，而是针对某一部分。因为，假使那些被撕裂之人的身体不得复活，但其他人的身体复活了，那么复活并没有被否定，只是某一类人除外而已。然而我诧异他们为何认为就连这方面也有疑问，似乎不是所有属土之物都归于土，并再次化为尘土。而且大海本身多半会将所吞没的人体冲到邻近的海岸上。即便不是如此，我想我们应当相信：将分散的结合、将散落的聚集，对于天主并非难事；宇宙服从他，沉默的元素顺从他，自然为他服务；仿佛赋予泥土生命比将其结合更为奇妙似的。\par

在\Place{阿拉伯}地区有一种鸟被称为\OriginalTerm{凤凰}{Phoenix}，它死后以新生的血肉汁液恢复，重新活过来：难道我们相信唯独人不会再复活吗？然而，我们通过众所周知的消息和文献权威得知，这种鸟的固定寿命为五百年，当它藉着本性的某种预兆知道自己生命终期临近时，它会为自己预备一个由乳香、没药和其他香料制成的匣子，当它的工作与时间同时结束时，它就进入匣子死去。然后从它的汁液中生出一只虫，逐渐成长为此鸟的模样，恢复其昔日的习性，借着双翼的划动，它再次开始其更新生命的历程，并偿还感恩的债。因为它从\Place{埃塞俄比亚}将那个匣子——或是它身体的坟墓，或是它复活的摇篮，在其中它离开生命而死，死了又复活——运到\Place{吕考尼亚}；因此，通过这只鸟的复活，那些地区的人便明白又一个五百年周期完成了。因此，对那只鸟来说，第五百年是复活之年，但对我们来说则是第一千年：它的复活在现世，我们的复活在世界的终末。许多人还认为这种鸟会自行点燃它的火葬柴堆，并从自己的灰烬中复生。\par

然而，若更深地探究本性，也许将为我们信仰提供一个更深层的理由：让我们的思绪回到人类受造的起源和开端吧。你们是男人和女人，你们并非不知人性相关之事，若你们中有人缺乏此知识，你们知道我们生于虚无。然而，我们如此伟大的存在，起源却是何等微小！若我不说得更明白，你们也明白我的意思，或说我不想明说之事。那么，这头，这奇妙的面容，其创造者我们无从得见，从何而来？我们看到这作品，它被造成以服务于各种目的和功用。这直立的身形、这高大的身躯、这行动的能力、这敏锐的感知、这直立行走的能力，从何而来？无疑，本性的器官我们并不知晓，但它们所产生的效果却为我们所知。你本人也曾是一粒种子，而你的身体正是那将要复活者的种子。听\Person{保禄}的话，你就会明白你就是这种子：\BibleQuote{播种的是可朽坏的，复活的是不可朽坏的；播种的是卑贱的，复活的是光荣的；播种的是软弱的，复活的是强健的；所播种的是属生灵的身体，复活的是属神的身体。} \BibleRef{格前 15:42-44} 那么，你也与其它东西一样被播种，如果你也将像其它东西一样复活，又何必诧异？但你相信那些事物，因为你看见了它们；你不相信这个，因为你没有看见：\BibleQuote{那些没有看见而相信的，才是有福的。} \BibleRef{若 20:29}\par

然而，在季节来临之前，那些事也不被人相信，因为并非所有季节都适合种子的生长。小麦在某一时播种，在另一时生长；葡萄在某时栽植，在另一时发芽的嫩枝开始抽条，枝叶茂盛，果实成形；橄榄在某时种植，在另一时，它仿佛孕育并满载着果实累累的子粒，因丰收而弯下了腰。但在各自的时节到来之前，产果受到限制，而结果之物并不由自己控制结果的时期。人们可以看到万物之母有时因霜霉而变形，有时光秃无产，有时绿意盎然、繁花似锦，有时干枯殆尽。任何地方若想永远披着外衣，从不脱下它种子的金色盛装，或草地的绿装，它在自身内便是荒芜的，且无法享受本可转移给其它产物的收益。\par

因此，即便你们既不愿凭信德、也不愿凭榜样相信我们的复活，你们也会凭经验而相信。因为许多产物，如葡萄树、橄榄树和各种果实，一年的终结乃是成熟的恰当时节；同样，对我们而言，世界的终结，犹如一年的终结，已为复活设定了恰当时机。死者的复活恰逢世界终结是合宜的，免得我们复活之后又须堕入此邪恶世代。为此，\Person{基督}受难，为将我们从这邪恶的世界中解救出来；免得此世的诱惑再次颠覆我们，若我们复活是为了犯罪，那复活便对我们有害了。\par

因此，我们既有了复活之理，也有了复活之时：理在于，本性在其一切产物中保持一贯，唯独在人类的生成上不会失败；时在于，万物都在一年的终了产生。因为世界的各季构成一年。既然一天是一，那么一年为一又何足为奇？因为主曾在一天内雇工人到葡萄园工作，他说：\BibleQuote{你们为什么站在这里整天闲着？} \BibleRef{玛 20:6}\par

64. 万物起始的原因乃是种子。外邦人的宗徒曾说，人的身体是一粒种子\BibleRef{格前 15:43}。如此，在播种之后，复活所需的实体便依次而来。但即使没有实体，也没有原因，谁又能认为天主从祂所愿之处并以祂所愿的方式重新造人是困难的？谁曾吩咐世界从无物质无实体中生出？请看高天，注视大地。星辰之火从何而来？太阳的圆轮与光芒从何而来？月亮的球体从何而来？高山的峰顶、坚硬的岩石、茂密的丛林从何而来？弥漫的空气、封闭或流溢的水从何而来？然而，如果天主从虚无中创造了这一切（因为\BibleQuote{祂一发言，万有造成；祂一出命，各物生成}），既然我们看见那不存在的得以产生，为何还应惊奇那曾经存在的会复活呢？\par

65. 令人惊奇的是，他们虽不信复活，却出于仁慈的关怀，设法使人类不致灭亡，因而主张灵魂会传递并迁入其他身体，好使世界不致消逝。但让他们说说看，哪一样更为困难：是灵魂迁移，还是回归；是回到那属于自己的，还是寻求新的居所。\par

66. 然而，让那些未受教导的人去怀疑吧。至于我们这些读过法律、先知、宗徒及福音的人，怀疑是不许可的。因为谁读了以下的话还能怀疑呢：\BibleQuote{那时，凡名录在那书上的，你的人民都能获救；许多长眠于尘土中的人要醒起来：有的得永生，有的永远蒙羞受辱。贤明之士要发光有如穹苍的光辉；引导众人归于正义者，要永远发光如同星辰。}所以，他论及那些睡眠的余者，为使人明白死亡并非永恒，它像睡眠一样只经历一段时期，时候到了就脱去；他显示，死后生命的进程，比死前在哀伤痛苦中度过的生命更美好，因为前者被比作星辰，后者被归于劳苦。\par

67. 还要我搜集别处所记吗？经上说：\BibleQuote{祢将使我复活，我必称谢祢。}还有那段圣\Person{约伯}在经历了此生苦难，并以德行的忍耐克胜一切逆境后，在复活中为现时的恶苦许下报酬的经文，他说：\BibleQuote{祢将使我这饱受苦难的身体复活。}\Person{依撒意亚}也向人民宣告复活，他说自己是主信息的传报者，因为我们读到：\BibleQuote{因为主的口已说话，到那日他们必这样说。}\BibleRef{依 25:8} 主的口所宣告人民应说的话，在后文有所阐述，那里记载着：\BibleQuote{上主，因敬畏祢，我们怀孕，并生出了祢救恩的精神，祢将这精神倾注于大地之上。地上的居民必将跌倒，坟墓中的人必将起来。因为从祢而来的甘露是他们的健康，但恶人的地必将灭亡。去吧，我的百姓，进入你的内室，隐藏片时，等待主的愤怒过去。}\par

68. 他多么恰当地以\Quote{内室}指代死者的坟墓，我们在其中隐藏片时，好能更胜于接受天主的审判，那审判将以我们罪有应得的义怒审判我们。那隐藏安息的人仍是活着的，仿佛从我们中间退避隐退，以免世间的愁苦以更紧的罗网缠住他；天上的圣谕藉众先知之声肯定，复活的喜乐为他存留，他们获释身体的健全将由天主的大能造成。甘露被用作标记也十分恰当，因为藉着甘露，大地上一切有生命的种子得以生长。那么，若我们衰败身体的尘土与灰烬，也因天上甘露的丰饶而充满活力，借这生命的滋润，肢体的形状得以重塑并重新连接，又有什么可奇怪的呢？\par

69. 圣先知\Person{厄则克耳}充分阐明并描述，枯骨如何重获活力，感觉恢复，动作加添，筋络归回，人体关节变得强壮；极干的枯骨如何被复原的肉所包裹，脉管的线路与血液的流动，被覆于其上的皮肤帷幔遮蔽。我们读到，复活的人体群象，似乎在先知的话音下涌现，可以看见在广阔的平原上，新的种子抽出新苗。\par

70. 古时的智者相信，种下龙牙，便在\Place{底比斯}地区长出一批武装的人；然而，已经确定，一种种子不能变为另一种植物，也不能产出与自身种子不同的果实，人不能从蛇生出，肉也不能从牙齿生出；那么，更要相信，所种的一切都按其本性复活，农作物不与种子相异，柔软的不从坚硬生出，坚硬的不从柔软生出，毒物不会变为血液；而是肉从肉再生，骨从骨再生，血从血再生，体液的各液从体液再生。你们这些异教徒，你们既能主张一种变化，怎能否认本性的恢复呢？你们既相信虚妄的寓言，岂能拒绝相信天主的圣谕、福音和先知呢？\par

71. 但让我们现在聆听先知本人的话语，他如此说道：\BibleQuote{主的手临于我，主的神带我出去，将我放在平原中，那里满是人的骨头；他领我在它们周围走了一遭，看，平原上骨头极多，而且非常干燥。他对我说：\InnerQuote{人子，这些骨头可以复生吗？}我答说：\InnerQuote{主，你知道。}他对我说：\InnerQuote{你向这些骨头讲预言，对它们说：干枯的骨头，听主的话吧！主对这些骨头如此说：看，我要使生命之神进入你们内，给你们放上筋，加上肉，包上皮，将我的神注入你们内，你们就复活了，并承认我是主。}我就按照他吩咐我的讲了预言。正当我讲这些预言时，看，发出一声巨响，大地震动。}\BibleRef{则 37:1}\par

72. 注意先知如何表明，在\OriginalTerm{生命之神}{Spirit of life}倾注于它们之上以前，骨头中已有了听觉和运动。因为，上面那干枯的骨头被吩咐聆听，仿佛它们有听觉；又，经先知的话指明，在此之后每块骨头都回到自己的关节，因为我们读到下文：\BibleQuote{骨头便互相结合起来，各归自己的关节。我又观看，看，筋和肉在它们上面出现，皮肤也从上面包了它们，但它们内还没有神。}\BibleRef{则 37:7-8}\par

73. 主的慈爱何等伟大，竟使先知成为未来复活的见证人，好让我们也能借他的眼目看见。因为并非人人都能被选作见证人，而是在那一位身上，我们所有人都成了见证人，因为谎言不会临于圣人，错谬也不会临于如此伟大的先知。\par

74. 骨头凭天主一声命令重新归于关节，这绝不应当显得难以置信，因为我们有无数的实例证明，自然曾顺从天上的命令；比如大地曾被命生出青草\BibleRef{创 1:11}，便生了出来；又比如岩石在棍杖触碰下为干渴的子民涌出水来；坚硬的石头因天主的仁慈，为被炎热灼伤的人倾泻了溪流。手杖变为蛇\BibleRef{出 4:3}，岂不正说明了，凭天主的旨意，无生命之物能产生有生命之物？你难道认为，骨头应命重新组合，比河水倒流或大海奔逃更难以置信吗？先知确实这样证实：\BibleQuote{大海看见而逃奔，约旦河也被迫倒流。}此处无可置疑，它由两个民族中一个获救、一个覆灭所证实：海浪静止而立，同时一面环绕一个民族，一面倒流向另一民族致其死命，乃是为了淹没后者，保全前者。我们又在福音本身中发现什么呢？主亲自在那里证明了，海洋因一句话而平息，乌云驱散，狂风止息，在平静的岸边，无言的元素顺从了天主。\par

75. 但让我们继续考察其他要点，好能看出，死者如何借着\OriginalTerm{生命之神}{Spirit of life}而复活，躺在坟墓中的如何起来，坟墓如何敞开：\BibleQuote{他对我说：\InnerQuote{人子，你应讲预言，对神讲说：主天主这样说：气息，你应四面八方来，吹在这些被杀的人身上，使他们复活。}我就照他命我的讲了预言，生命之神进入了他们内，他们便复活了，并且站了起来，实在是一支极庞大的军队。主对我说：\InnerQuote{人子，这些骨头就是以色列全家族。他们常说：我们的骨头干枯了，我们的希望消失了，我们全完蛋了。因此，你应讲预言，对他们说：主天主这样说：看，我要亲自打开你们的坟墓，将你们从坟墓里领出来，回到以色列地域。当我打开你们的坟墓，将我的民族从你们的坟墓里领出来的时候，你们便知道我是主。我要把我的神注入你们内，使你们复生，叫你们安居在自己的地域，那时，你们便知道我是主；我言出必行——主的断语。}}\BibleRef{则 37:9}\par

76. 我们在此注意到，\OriginalTerm{生命之神}{Spirit of life}的运作再次恢复；我们明白了，坟墓敞开时死者以何种方式复活。当真不可思议吗？死者的坟茔因主的命令而开启，当整个大地从极边之处被一声霹雳震动，海水漫过界限，又重新阻挡波涛的路线？最后，那相信死者将复活的人——\BibleQuote{在一霎时，眨眼之间，在\OriginalTerm{末次吹响号筒}{the last trump}时（因为号筒将要响），}\BibleRef{格前 15:52}——\BibleQuote{将首先被提到云中，在空中迎接基督；}\BibleRef{得前 4:17}那未曾信的人将被留下，并因自己的不信而定自己的罪。\par

77. 主也在福音中，现在让我们来看一些实例，向我们展示我们将以何种方式复活。\Quote{因为祂不只复活了拉匝禄，更复活了所有人的信德；你若相信，在你阅读时，你那曾死的灵魂也与拉匝禄一同复活。}主走到坟墓前，大声呼喊：\BibleQuote{拉匝禄，出来！}\BibleRef{若 11:43}除了要给我们一个明显的证据，为未来的复活树立一个榜样之外，还有什么意思呢？祂何以大声呼喊，仿佛祂不习惯于在圣神内工作、在沉默中命令，而只是为了显示经上所说的：\BibleQuote{在一霎时，眨眼之间，在\OriginalTerm{末次吹响号筒}{the last trump}时，死人要复活成为不朽的}？\BibleRef{格前 15:52}因为高声呼喊对应着号筒的响声。祂呼喊：\BibleQuote{拉匝禄，出来！}\BibleRef{若 11:43}为何加上名字？恐怕是为免得一人顶替另一人被复活，或是为使复活显为命定的，而非偶然的。\par

78. 这样，死者听见了，便从坟墓中出来，手脚缠着殓布，脸上包着汗巾。试想，你若能，他如何闭着眼睛行走，用被捆绑的双脚指引脚步，肢体被束缚却如同自由人一样行动。绷带留在他身上，却没有束缚他；他的眼睛被蒙住，却能看见。这样，他就看见了谁在复活，谁在行走，谁在离开坟墓。因为当天主的命令大能运行时，自然不需要自己的功能，被带到极限时，不再遵循自己的路径，而是遵循天主的旨意。死亡的捆绑在坟墓的捆绑之前就破裂了。行动的能力在行动的方法被提供之前就运用了。\par

79. 如果你对此感到惊奇，想想是谁发出了命令，你就会停止惊奇；\Person{耶稣基督}，天主的能力，生命，光明，死者的复活。那能力扶起了倒卧的人，那生命引领了他的脚步，那光明驱散了黑暗，恢复了他的视力，那复活更新了生命的恩赐。\par

80. 或许你会因犹太人搬开石头、解开殓布而感到不安，或许你会担心谁将搬开你坟墓的石头。好像那位能恢复神魂的，却不能移开石头；或那位能使被捆绑者行走的，却不能挣断束缚；或那位曾光照蒙蔽之眼的，却不能揭开面容；或那位能更新自然进程的，却不能劈开石头！但他们为了让那些不愿用心相信的人相信自己的眼睛，就搬开石头，看见尸体，闻到臭气，解开殓布。他们看见他复活，就不能否认他是已死的人；他们看见死亡的标记和生命的证据。如果他们在忙碌时，因这劳作本身而悔改呢？如果他们听见时，相信自己的耳朵呢？如果他们看见时，被自己的眼睛教导呢？如果他们解开束缚时，释放自己的心灵呢？如果当\Person{拉匝禄}被解开时，百姓得释放，当他们释放\Person{拉匝禄}时，自己归向主呢？最后，许多来到\Person{玛利亚}那里的人，看见所发生的事，就相信了。\par

81. 这并非我们的主\Person{耶稣基督}所行的唯一实例，祂还复活了其他人，好叫我们至少因更多的实例而相信。祂因那寡妇母亲的眼泪而感动，复活了那少年人，当祂上前按住棺架，说：\BibleQuote{少年人，我吩咐你，起来吧！那死人就坐起来，并开口说话。}\BibleRef{路 14:7} 他一听见，就立刻坐起来，立刻说话。因此，能力的运作是一回事，自然的秩序是另一回事。\par

82. 我又该怎样论及会堂长的女儿呢？她死时，众人哭泣，吹笛者奏乐；因确信死亡，正在举行葬礼仪式。主的话一出口，灵魂如何迅速回归，复苏的身体如何起来，又进食，好使人相信生命的证据！\BibleRef{谷 5:38}\par

83. 我们为何要惊奇于灵魂因天主的话而恢复，肉体归于骨骼，当我们记得先知身体的触碰能使死者复活时？\Person{厄里亚}祈祷，复活了死去的孩子。\Person{伯多禄}因\Person{基督}之名吩咐\Person{塔彼达}起来行走，\BibleRef{宗 9:40} 穷人们因她服事他们的食物而欢欣喜信，我们难道不该为了我们的救恩而相信吗？他们用眼泪换来了别人的复活，我们难道不相信\Person{基督}的苦难为我们换来的复活吗？当祂断气时，为了表明祂是为我们的复活而死，就完成了复活的进程；因为就在那时\BibleQuote{祂大声呼喊，断了气，大地震动，岩石崩裂，坟墓自开，许多已安眠的圣人的身体复活了，在祂复活后，他们出了坟墓，进入圣城，显现给许多人。}\BibleRef{玛 27:50}\par

84. 如果这些事在祂断气时发生，我们为何认为祂将来审判时再来时这些事不可信呢？尤其因为这早先的复活是那未来复活的保证，是那将来实现的预象；事实上，它本身就是真实，而非预象。那么，在主复活时，是谁打开了坟墓，向那些复活者伸手，指示他们前往圣城的道路？若没有别人，那必定是天主的能力在死者体内运行。人看见天主的工程，难道还要寻求人的帮助吗？\par

85. 天主的行动无需人的帮助。天主命令诸天存在，便成就了；祂决定创造大地，便创造了。谁用肩扛来石头？谁提供费用？谁在天主工作时为祂提供帮助？这一切都是片刻间造成的。你要知道有多迅速吗？\BibleQuote{祂一发言，它们便造成。}如果元素因一句话而产生，为何死者不能因一句话而复活？因为他们虽然死了，却曾生活过，曾拥有生命的气息以感知，拥有行动的力量；并且，从未能拥有生命与保持无生命之间，有着极大的区别。魔鬼说：\BibleQuote{命这块石头变成饼吧。}\BibleRef{路 4:3} 它承认因天主的命令自然能被转变，你不相信因天主的命令自然能被重塑吗？\par

86. 哲学家们就太阳的轨道与天体的系统争论不休，且有人以为应当相信他们，尽管他们对自己所谈论的一无所知。因为他们既未升入高天，也未测量过苍穹，更未用双眼观察过宇宙；因为他们中没有人在起初与天主同在，没有人论及天主说：\Quote{当祂预备高天时，我已与祂同在，我作为技师与祂同在，我是祂所喜悦的。}那么，如果人们相信他们，难道不相信天主吗？天主这样说：\BibleQuote{正如我造的新天新地，在我面前怎样存留——主这样说——你们的名字与后裔也将怎样存留；月复一月，安息日复安息日，一切有血肉的都要来到我面前，在\Place{耶路撒冷}朝拜——主天主这样说；他们将要出去，观看那些背叛我的人的尸体；因为他们的虫不死，他们的火不灭，他们将成为一切有血肉者的鉴戒。}\BibleRef{依 66:22}\par

87. 如果大地与苍天都将更新，我们为什么要怀疑那天地为他而造的人能被更新？如果\OriginalTerm{背命者}{transgressor}被保留以待惩罚，为什么义人不能为光荣而被保守？如果\OriginalTerm{罪虫}{worm of sins}不死，义人的血肉又怎能消逝？因为复活——该词的形式本身即表明——是指那已跌倒的再起来，那已死的再复生。\par

88. 这正是正义的规程与基础：由于灵魂与身体的行为为二者所共有（因为灵魂所构想的，身体已实施），二者都应接受审判，也都应或被交付惩罚，或被保留为光荣。否则，似乎极不协调：既然心神的法律与肉欲的法律交战，且心神常常——当寓居在人内的罪恶活动时——行其所恨恶的事；那么，因与他人共有的过错而有罪的心神，竟要受罚，而作为恶之肇始的血肉，竟得享安息；而且，那并非单独犯罪的，竟要单独受苦，或者，那并未单独在恩宠助佑下作战的，竟要单独获得光荣。\par

89. 这理由，除非我错了，是完备且正义的，但我并不向基督寻求理由。如果我被理由说服，我就拒绝了信德。\Person{亚巴郎}信了天主，\BibleRef{创 15:6}我们也该信祂，好使我们这些他种族的继承人，也成为他信德的继承人。\Person{达味}同样信了，所以他才说话；我们也该信，好使我们能说话，因为知道：\BibleQuote{那使主耶稣复活者，也必要使我们与耶稣一同复活。}\BibleRef{格后 4:14}因为从不撒谎的天主曾许下此事；\OriginalTerm{真理}{Truth}在祂的福音中也如此许下，当祂说：\BibleQuote{这是派遣我来者的旨意：凡祂交给我的，叫我连一个也不失落，且在末日使他复活。}\BibleRef{若 6:39}祂认为只说一次不足够，更以明确的重复加以强调，因为紧接着说：\BibleQuote{因为这是我父的旨意，即派遣我来者的旨意：凡看见子而信从祂的人，必获得永生，并且在末日我要使他复活。}\par

90. 说这话的是谁？正是那位死了却使许多亡者的身体复活者。如果我们不信天主，难道我们也不信证据吗？既然祂甚至做了祂未曾应许的事，难道我们不信祂所应许的吗？祂有什么理由去死，若祂没有理由复活呢？因为，既然天主不能死，\OriginalTerm{智慧}{Wisdom}也不能死；而既然未曾死的不能复活，祂便取了能死的\OriginalTerm{血肉}{flesh}，好使那按其本性会死的死去，而那死过的得以复活。因为复活除非藉着人，便无法实现；因为，\BibleQuote{就如死亡因一人而来，死者的复活也因一人而来。}\BibleRef{格前 15:21}\par

91. 所以，人复生，因为人死了；人复活了，但由天主使他复活。那时，他是按血肉的人，而今，天主是\OriginalTerm{万有中的万有}{all in all}。\BibleRef{格前 15:28}因为如今我们不再按血肉认识基督了，\BibleRef{格后 5:16}但我们享有那血肉的恩宠，好使我们认识祂为安眠者的\OriginalTerm{初果}{firstfruits}，\BibleRef{格前 15:23}死者的\OriginalTerm{首生者}{firstborn}。\BibleRef{哥 1:18}初果无疑与其余的果实具有同一性质和种类，首先献上的这些是为祈求更丰盛的收成，作为一切恩赐的圣洁感恩祭，并作为那已恢复的本性的一种奠祭。因此，基督是安眠者的初果。但这仅是那些属祂而安眠者——即那些仿佛从死亡中解脱，被一种甜蜜的沉睡所掌握的人——的初果，还是所有死者的初果呢？\BibleQuote{就如在基督内，众人都死了，照样，在基督内，众人都要复活。}\BibleRef{格前 15:22}所以，正如死亡的初果在\Person{亚当}内，同样，复活的初果在基督内。\par

92. 所有人都要复活，但不要有人灰心，也不要让义人因复活的共同命运而悲伤，因为他等候他德行的卓越果实。所有的人的确都要复活，\BibleRef{格前 15:23}但正如宗徒所说：\BibleQuote{各人依照自己的\OriginalTerm{次序}{order}。}\OriginalTerm{天主的仁慈}{Divine Mercy}的果实为所有人共有，但\OriginalTerm{功绩}{merit}的顺序却不同。白昼光照众人，太阳温暖众人，雨水以柔和的阵雨滋润所有人的田地。\par

93. 我们所有人都出生了，我们也必将全都复活；但在每一种状态中，无论是生还是复活后的生活，恩宠不同，状况也不同。因为，\BibleQuote{在顷刻眨眼之间，在末次吹号筒时，死人必要复活成为不朽的，我们也必要改变。} \BibleRef{格前 15:52} 而且，在死亡本身中，有些人安眠，有些人活着。安眠是好的，但生命更好。所以宗徒唤醒那安眠的人进入生命，说：\BibleQuote{你这睡眠的，醒来吧！从死者中起来吧！基督必要光照你。} \BibleRef{弗 5:14} 因此他受唤醒，为能生活，为能像保禄一样，为能说：\BibleQuote{我们这些活着存留到主来临时的人，决不会在已安眠的人以前。} \BibleRef{得前 4:14} 他在这里不是谈论我们共同的生活方式和我们同样享有的气息，而是谈论复活的功劳。因为，在说了\BibleQuote{那些在基督内死去的人要先复活}之后，他进一步说：\BibleQuote{然后我们这些活着还存留的人，同时与他们一起要被提到云彩上，到空中迎接主。} \BibleRef{得前 4:17}\par

94. 保禄确已死去，并因他的光荣受难，以肉身的生命换取了永久的荣耀；那么当他写道他将被活生生提到云中与基督相遇时，他是在欺骗自己吗？我们也在\BibleBook{创}{世纪}中读到哈诺客的事\BibleRef{创 5:24}，以及厄里亚的事，你也必将在圣神内被提到。看，厄里亚的火车，看，那火，虽未看见却已备妥，为使义人得以升天，无辜者被带去，而你的生命不知死亡。因为宗徒们确实不知死亡，正如经上所说：\BibleQuote{我实在告诉你们：站在这里的人中，就有些人在未尝到死味以前，必要看见人子来到自己的国内。} \BibleRef{玛 16:28} 因为他活着，在他内没有任何可死之物，他没有任何来自埃及的鞋或束缚，而是在放下这身体的服役之前就已脱去了。所以，不仅是哈诺客一人活着，因为不是他一人被接去；保禄也被提去与基督相会。\par

95. 圣祖们也活着，因为天主若不能被称作亚巴郎、依撒格和雅各伯的天主，除非死者是活着的；因为祂不是死人的天主，而是活人的天主。 \BibleRef{路 20:38} 而我们，若愿意效法前人的行为和习惯，也必将活着。我们惊叹圣祖们的赏报，让我们效法他们的忠信；我们讲述他们的恩宠，让我们追随他们的服从；让我们不要被贪欲引诱，陷入世界的罗网。让我们抓住机会，把握法律的诫命，我们圣召的仁慈，受苦的渴望。圣祖们离开了自己的土地，让我们在决意上从肉身的权势中离去；让我们在决意上像他们一样流亡；但他们不把那由敬畏天主而导致的视为流亡，也非迫不得已。他们以另一片土地更换了自己的故土，让我们以天国更换尘世；他们更换了地上的居所，让我们在精神中更换。智慧曾向他们显示那布满星辰的天空， \BibleRef{创 15:5} 让智慧光照我们心灵的眼目。这样，预像与真理相符，真理也与预像相符。\par

96. 亚巴郎，乐于款待旅客，对天主忠信，事奉殷勤，服务迅速，在预像中看见了天主圣三；\BibleRef{创 18:2} 他在款待之上加上了宗教的义务，当他看见三位时，却朝拜一位；他保持了位格的区别，却称呼一位主；他承认一个能力，却向三位献上礼物的尊荣。在他内说话的不是学识，而是恩宠；他比我们这些已学习的人，更能相信他所未学习的事。没有人曾伪造过真理的呈现，所以他看见三位，却朝拜一体。他拿出三斗细面，并宰杀一头牺牲，认为一次祭献已足够，但却是三重的礼物；一头牺牲，却是三位的奉献。\EditorialNote{创 14} 至于那四王，谁不理解为，他以此征服了物质受造界的元素以及一切属地的事物，作为一种预像，预示了主的受难？他在战争中忠信，在胜利中节制，因为他宁愿不因人的馈赠致富，而是因天主的恩赐致富。\par

97. 他相信自己在老年时还能生子， \BibleRef{创 15:6} 并且当他成为父亲时，断定自己能献上儿子为祭；当责任扶助老翁的右手时，他的父爱并未战栗， \BibleRef{创 22:11} 因为他知道，被祭献的儿子比完整的儿子更蒙天主悦纳。因此，他带着他的爱子去献为祭，并迅速献上了他迟来得到的儿子；当爱子称他为\Quote{父亲}，他回答\Quote{我儿}时，他并不因被称为父亲而受阻。这些称呼是爱的珍贵保证，但天主的诫命更受爱戴。所以，尽管他们心心相印，他们的主意却坚定不移。父亲的手将刀举过儿子的身体，父亲的心挥下一击，为使定案不致落空；他生怕刀落不准，生怕右手乏力。他感到了父爱的涌动，却不从服从的工作上退缩，反而加速他的服从，即便当他听到来自天上的声音时亦是如此。那么，让我们将天主置于我们所爱的一切之上，父亲、兄弟、母亲，好叫天主保全我们所爱的人，正如在亚巴郎身上，我们看到的多是慷慨的赏报者，而非仆人。\par

\Person{亚巴郎}这位父亲确曾献上他的儿子，但天主不以流血而悦纳孝爱的服从。祂在灌木丛中指给他一只公绵羊\BibleRef{创 22:13}，替代那孩子，为使儿子得以归还父亲，而祭品并不缺少司祭。这样，\Person{亚巴郎}未因儿子的血而受玷污，天主也未失去祭献。先知发言了，既未流于自夸，也未固执到底，而是以公绵羊换取了孩子。由此更显明他如何虔诚地献上，如今又如此乐意地将他领回。而你，你若将你的礼物献给天主，并不会失去它。但我们却紧守己物；天主为我们赐下了祂的独生子，\BibleRef{罗 8:32}而我们拒绝献出我们的。\Person{亚巴郎}看见了这事，认出了这奥迹：救恩要从那树临于我们，他也没有忽略，在同一祭献中，看似被献上的是一位，真能受死的却是另一位。\par

那么，让我们效法\Person{亚巴郎}的虔诚，效法\Person{依撒格}的良善，效法他的贞洁。这人显然良善而贞洁，对天主满怀虔诚，对妻子保持贞洁。他不以恶报恶，向驱逐他的人让步，待他们悔改后又重新接纳他们，对傲慢并不暴躁，对善意也不固执。当他离开他人时远避纷争，当他重新接纳他们时乐于宽恕，当他宽恕他们时恩慈更为丰厚。人们寻求与他交往，他还另行摆设欢乐的筵席。\par

在\Person{雅各伯}身上，也让我们效法\Person{基督}的预像，愿我们的行为中带有些许他的样式。我们若效法他，便将与他有份。他服从母亲，谦让兄长，服事岳父，从羊群的增加而非分割中寻求工价。在所得如此丰厚的部分里，没有贪婪的分割。那记号也并非无目的：由地及天的梯子，\BibleRef{创 28:12}其中预示了将来藉\Person{基督}的十字架，人与天使之间的交往，他的大腿瘸了，\BibleRef{创 32:25}好让他在这腿上认出自己肉身的后嗣，并藉大腿的瘸预示后嗣的受难。\par

那么，我们看到，天堂为德行敞开，而且这不只是少数人的特恩：\BibleQuote{将有许多人从东方和西方来，从北方和南方来，在天主的国里坐席。}\BibleRef{玛 8:11}这表达了永享安息的喜乐，因为灵魂的躁动止息了。让我们在行为上效法\Person{亚巴郎}，好使他接我们到他的怀抱，以慈爱的拥抱抚慰我们，如同\Person{拉匝禄}——他谦卑的后嗣，被自己的特有德行所环绕。这蒙天主悦纳的圣祖的追随者，不是以肉身的怀抱，而是以犹如善工的外衣抚慰我们。\BibleQuote{你们不要自欺，}\Person{宗徒}说，\BibleQuote{天主是轻慢不得的。}\BibleRef{迦 6:7}\par

我们已看到，不信复活是多么重大的冒犯；因为我们若不复活，\Person{基督}就白白死了，\Person{基督}也没有复活。\BibleRef{格前 15:13}因为祂若不为我们复活，则祂根本未曾复活，因祂无需为自己复活。宇宙在祂内复活了，天在祂内复活了，地在祂内复活了，因为将有新天新地。\BibleRef{默 21:1}但对于死亡之权所不能拘禁的那一位，复活又有何必要？因为祂虽作为人而死，却在阴府中也是自由的。\par

你想知道祂如何自由吗？\Quote{我成了一个无助的人，在死者中自由。}祂被称为自由者，确然恰当，因祂有能力使自己复活，正如经上所载：\BibleQuote{你们拆毁这座圣殿，三天之内我要把它重建起来。}\BibleRef{若 2:19}祂被称为自由者，也确然恰当，因祂曾下降，为拯救他人。因为祂成了人，并非只是外表，而是真实地如此造成，因为祂是人，谁又能认识祂呢？因为，\BibleQuote{祂取了奴仆的形体，与人相似，形状也一见如人；祂自谦自卑，听命至死，}\BibleRef{斐 2:7}为使我们藉那听命看见祂的光荣，就是\BibleQuote{圣父独生子的光荣}，\BibleRef{若 1:14}正如圣\Person{若望}所说。这样，圣经的叙述得以保持，若在\Person{基督}身上，独生子的光荣与完美的人性都保持不失。\par

因此，祂无需帮手。因为祂创造世界时无需帮手，所以救赎世界时也无需帮手。没有使者，没有信差，而是主亲自使之完备。\Quote{祂一发言，万有造成。}主亲自使之完备，在每一部分都亲自为之，因为万有都藉祂而有。那位在祂内万有受造、靠祂而万有存在的主，谁能帮助祂？\BibleRef{哥 1:17}那位在一瞬间造作万有、并在末次号角时复活死者的主，谁能帮助祂？\BibleRef{格前 15:52}所谓\Quote{末次}，并非祂不能在第一次、第二次或第三次就使他们复活，而是遵循一种次序，并非为了最终克服困难，而是为了完成预定的数目。\par

但我以为，如今正是谈论号角的时候了，因为我的演讲已近尾声，愿号角也作为我发言结束的标志。我们在\Person{若望}的《\WorkTitle{默示录}》中读到七枝号角，由七位天使领受。\BibleRef{默 8:2}在那里你又读到，当第七位天使吹响号角时，有大声自天而降，说：\BibleQuote{世上的国已成了我们的天主和祂的基督的国，祂要为王，至于无穷之世！}\BibleRef{默 11:15}号角一词也用来指声音，正如你读到：\BibleQuote{此后，我看见天上有门开了；我初次听见的那好像号角的声音，对我说：\InnerQuote{你上到这里来，我要把以后必定发生的事指示给你。}}\BibleRef{默 4:1}我们也读到：\Quote{在月朔之日，你们当吹号角}；另外又读到：\Quote{当以号角之声赞美祂。}\par

106. 因此，我们应当尽一切力量留意号角的意义，免得我们像老妇人那样，把它们仅当作故事的一部分，如果我们认为这些事物配不上属灵的教导，或有损圣经的尊严，便会陷入危险。因为当我们读到，我们的战斗不是对抗血和肉，而是对抗那些在高天领域的属灵邪恶势力\BibleRef{弗 6:12}，我们就不应想到属血肉的武器，而应想到那些在天主面前大有能力的武器\BibleRef{格后 10:4}。一个人看见号角或听见它的声音是不够的，除非他明白这声音的意义。因为如果号角发出了模糊的声音，谁还会准备作战呢？\BibleRef{格前 14:8}因此，我们理解号角之声的含义至关重要，免得我们在听见或发出这类号角声时，显得像是化外之人。所以，当我们发言时，让我们祈求圣神为我们解释这些意义。\par

107. 那么，让我们探究我们在\WorkTitle{旧约}中所读到的关于各种号角的记载，考虑到法律为犹太人所规定的那些节庆，乃是高天之欢乐与天上节庆的影子。因为这里是影子，那里是真实。让我们努力通过影子来达到真实。这真实的预像以这样的方式表达出来，我们在经上读到主对梅瑟说：\BibleQuote{你去向以色列子民说：七月，月初，你们应完全休息，应吹号角为纪念，应召开圣会；任何劳动工作都不可做，并应向主献全燔祭。} \BibleRef{肋 23:24} 在\BibleBook{户籍}{纪}中：\BibleQuote{主训示梅瑟说：你应制造两个银号角，用锤子打成，为召集会众，为扎营时号令起营之用。当你吹号角时，全会众都应聚集在会幕门口。如果你只吹一个号角，以色列的众领袖和首长都应聚集到你这里；你第一次吹号角发出信号时，扎在东方的营都要起程。你第二次吹号角发出信号时，他们拔营前往，安营朝向黎巴嫩方向。你第三次吹号角发出信号时，他们拔营前往，安营朝向北方\OriginalTerm{北方}{Boream}。你第四次吹号角发出信号时，他们拔营前往，安营朝向北方\OriginalTerm{北方}{Aquilonem}。他们拔营前行时，应吹号角发出信号。但在召集会众时，你只吹号角，不发出信号。亚郎的子孙作司祭的应吹号角：这是你们世世代代永久的法令。当你们在本国内，出外作战，抵抗来侵略的仇敌时，你们应吹号角发出信号，你们便被主你们的天主记念，从你们的敌人手中被救出来。此外，在你们的庆日、节日、月朔，献全燔祭及和平祭时，你们也应吹号角，如此你们的祭献在天主面前将成为你们的纪念：我是主你们的天主。} \BibleRef{户 10:1}\par

108. 那么，怎样呢？难道我们只是凭吃喝来评判节日吗？但不要让人在吃喝上论断我们；\BibleQuote{因为我们知道，法律是属灵的。} \BibleRef{罗 7:14} \BibleQuote{所以不要让人在饮食上，或在节期、月朔或安息日等事上论断你们，这些都不过是未来事物的影子，但实体却是基督。} \BibleRef{哥 2:16} 那么，让我们寻求基督的身体，就是圣父从天上发出的声音，如同末次的号角，在犹太人说打雷了的时候，向你们所显明的；\BibleRef{若 12:29} 这基督的身体，末次的号角也要再次彰显出来；因为\BibleQuote{主要亲自由天降下，伴随总领天使的呼声和天主的号角声，那些在基督内死了的人要复活起来；} \BibleRef{得前 4:16} 因为\BibleQuote{尸首在哪里，鹰也在那里。} \BibleRef{路 17:37} 基督的身体在哪里，真实就在哪里。\par

108. 那么，第七支号角似乎象征那周的安息日，它不仅以日、年及周期计算（为此喜年之数也是神圣的），也包括第七十年，那时百姓在流徙七十年后返回了耶路撒冷。在这神圣数字的遵守上也绝未忽略百计、千计，因为主并非无意地说：\BibleQuote{我为我自己留下了七千人，这些人的膝从未向巴耳屈过。} 因此，未来安息的影子在此世的日、月、年中得以预表，为此，梅瑟命令以色列子民在第七月，月的第一日，借着\Quote{吹号角之纪念}为全体立定安息，不可做任何奴役之工，而应向天主献祭，因为在这周之末，就是世界的安息日，要求我们的，是属灵而非属肉身的工作。因为那属肉身的是奴役的，因为肉身服务于灵魂，而无罪使人自由，罪过使人沦为奴隶。\par

所以，属神的事必须如同在镜中和谜语中那样被认识；\BibleQuote{因为我们现在是对着镜子观看，模糊不清，到那时就要面对面的观看了。}\BibleRef{格前 13:12} 如今我们按血肉争战，那时便要在\Person{圣神}内得见神圣的奥秘。那么，让真法律的特质彰显在我们这些按\Person{天主}的肖像行事的人的生活方式中，因为法律的\OriginalTerm{影子}{shadow}现在已经消逝了。肉性的犹太人拥有影子，\OriginalTerm{肖像}{image}则属于我们，\OriginalTerm{实体}{reality}则属于那些将要复活的人。因为我们知道，依法律而言，有这三者：影子、肖像或\OriginalTerm{模样}{likeness}，以及实体；影子在法律中，肖像在福音中，真理在审判中。但一切是\Person{基督}的，一切都在\Person{基督}内；如今我们不能按实体看见\Person{基督}，但我们看见\Person{基督}，如同在一种未来事物的模样中，这未来事物的影子我们已在法律中见过了。所以，\Person{基督}不是\Person{天主}的影子，而是肖像，不是空洞的肖像，而是实体。因此，法律是藉\Person{梅瑟}颁赐的，因为影子是藉着人，模样是藉着法律，实体是藉着\Person{耶稣}。因为实体不能来自实体以外的任何源头。\par

那么，如果有人渴望看见这天主的肖像，就必须爱\Person{天主}，好能蒙\Person{天主}所爱；不再作仆人，而是作朋友，因为他遵守了\Person{天主}的诫命，好能进入\Person{天主}所在的云中。\BibleRef{出 24:15} 让他为自己制造两支\OriginalTerm{属理智的}{reasonable}号角，用经过锤炼、证明的银子做成，也就是说，由珍贵的话语组成并加以修饰，从中发出的不是尖锐刺耳、令人恐惧的声音，而是对\Person{天主}的崇高感谢，随不停的欢跃而倾流。因为借由这种号角的声音，死人得以复活，确实不是凭借金属的声音，而是被真理之言唤醒。或许那两支号角就是\Person{保禄}藉着\Person{圣神}说话时所用的，他说：\BibleQuote{我要以神魂祈祷，也要以理智祈祷；我要以神魂歌咏，也要以理智歌咏；}\BibleRef{格前 14:15} 因为两者缺一，似乎决不能有全备的说话。\par

然而，并非人人都能吹每支号角，也不是人人都有权召集全体会众，这特权只授予\Person{司铎}，\BibleRef{户 10:8} 以及吹响号角的\Person{天主}的仆役，好叫凡听见，并追随到上主荣耀所在之处，及早决心来到\Place{约柜的帐幕}的人，也能得见天主的工程，并为他全部后裔的延续而堪当承受那指定而永恒的居所。因为，当\Person{圣神}的恩宠与灵魂的力量协同行动时，战争便告终结，敌人也被驱散。\par

而且，如果一个人心里相信，口里承认，这些号角也就成为有益的号角；\BibleQuote{因为心里相信，可使人成义，口里承认，可使人获得救恩。} 因为凭借这双重的号角，人得以抵达那圣地，就是复活的恩宠。所以，让它们时常向你们吹响，使你们能时常听见\Person{天主}的声音；愿天使和先知的话语时常激励和推动你们，使你们能向往天上的事。\par

\Person{达味}心中正是默想着这个目的，当他说道：\BibleQuote{我要进入奇妙帐幕的所在，直到\Person{天主}的殿里，随着欢跃和称谢的歌声，赴宴者的喧嚷。} 因为，敌人不仅被这些号角的声音所胜；若无它们，也没有欢庆、节日或\OriginalTerm{月朔}{new moons}。因为没有人能守节日或月朔，除非他领受了\OriginalTerm{圣言}{Divine Word}的应许，并相信由此而来的信息；在那些节期里，他渴望让自己摆脱肉体的享乐和世俗的操劳，充满\Person{基督}的光明。而且，祭献本身，若没有口里的\OriginalTerm{宣认}{confession}相伴，也不能蒙\Person{基督}悦纳；按照惯例，口里的宣认在司铎奉献祭品时，激励民众恳求\Person{天主}的恩宠。\par

所以，让我们成为上主的宣讲者，在号角声中赞美\Person{天主}，不要轻看或小觑这号角的力量，而要发出能充满心灵之耳、深入我们最深处意识的话语，使我们不至于以为那适于肉身的事也可应用于\OriginalTerm{天主性}{Godhead}，也不要以人的能力衡量\Person{天主}大能的伟大，以至去探究人如何能复活，或以何种身体再来，或那些已分解的如何能再结合，已失去的如何能恢复，因为这一切，只要\Person{天主}定意，就立刻成就。而且，我们所期待的并不是一种肉身感官能分辨的号角声，而是天上尊威那不可见的能力在运作；因为对\Person{天主}而言，愿即是行；我们无须探究复活所需的力量，却要为自己寻求复活的果实。这一切将更易成就，如果我们摆脱过犯，达到\OriginalTerm{属神奥秘}{spiritual mystery}的圆满，那更新了的肉身从\Person{圣神}领受恩宠，灵魂从\Person{基督}获得永恒光明的辉煌。\par

然而，这些奥秘不仅关乎个人，也关乎全人类。试看依照\OriginalTerm{法律的预像}{type of the Law}，\OriginalTerm{恩宠的秩序}{order of grace}。当第一支号角吹响时，它聚集朝向东方的人，即那些为首的、被选的人；第二支吹响时，聚集那些功绩几乎相等的人，他们位向\Place{黎巴嫩}，已抛弃了万民的愚妄；第三支吹响时，聚集那些仿佛在这世界的海洋上颠簸，被今生的波浪冲来冲去的人；第四支吹响时，聚集那些绝未能用\OriginalTerm{属神话语}{spiritual utterance}的命令充分软化自己的心硬，因而被说成是朝向北方的——因为，按\Person{撒罗满}所说，北方是凛冽的风。\par

116. 所以，虽然所有人都在一瞬间复活，但每个人都是按自己的功绩次序复活的。因此，那些首先顺服热心的催促，如同在信德的曙光之前出发的人，首先领受了永恒太阳的光芒。这可以正确地论及旧约的圣祖们，或福音之下的宗徒们。第二类是那些弃绝外邦人的仪式，从不洁的错误中出来接受教会训练的人。所以，前者是圣祖们，后者是外邦人，因为信德之光由前者开始，在后者中存留直到世界的终结。第三和第四位复活的，是那些在南方和北方的人。大地分为这四部分，一年由这四部分组成，地由这四部分完成，教会由这四部分聚集而成。因为所有被视为加入圣教会的人，因呼求天主的名，都将获得复活的殊荣和永恒福乐的恩宠，因为\BibleQuote{他们将从东从西，从北从南而来，坐在天主的国里。}\BibleRef{路 13:29}\par

117. 因为基督用来环绕祂世界的并非小光：因为\BibleQuote{祂的出发是从天的高处，祂的行程至天的高处，没有谁能隐藏自己免受祂的热力。}\BibleRef{咏 19:7}因为祂以祂的良善光照众人，不愿拒绝愚昧的人，而愿改正他们；不愿将心硬者排除在教会之外，而愿软化他们。因此，教会在\WorkTitle{雅歌}中、基督在\WorkTitle{福音}中邀请他们说：\BibleQuote{凡劳苦和负重担的，你们都到我跟前来，我要使你们安息。你们背起我的轭，跟我学吧！因为我是良善心谦的。}\BibleRef{玛 11:28-29}\par

118. 你也可以认出教会邀请的声音，因为她曾说：\BibleQuote{北风啊，醒来吧！南风啊，吹来！吹拂我的园子，使我的香液流出来。让我的良人来到他的园中，吃他珍奇树的果实。}\BibleRef{歌 4:16}因为，啊，圣教会，你那时就已知道，从这些人身上也会为你产生丰硕的成果，你向你的基督许诺从这样的人得到果实。你首先说过，你被引入君王的寝室，爱慕祂的胸怀胜似美酒，因为你爱了那爱你的，寻求了那养育你的，并为宗教的缘故轻视了危险。\par

119. 然后，啊，净配，你要从\Place{黎巴嫩}被召来，在主的审判中，你全然美丽，毫无瑕疵。因为经上记载说：\BibleQuote{我的爱卿，你全然美丽，毫无瑕疵。我的新娘，从\Place{黎巴嫩}下来，从\Place{黎巴嫩}下来吧！}\BibleRef{歌 4:7}\par

120. 随后，你既不惧怕汹涌的水，也不惧怕从\Place{黎巴嫩}奔流而下的溪流，你呼唤北风和南风，希望它们吹拂你的园子，使你的香液流到别人身上，好让你能将他人的丰硕成果作为你丰收的献礼，呈给基督。\par

121. 因此，\BibleQuote{凡遵守这预言的话的，是有福的。}\BibleRef{默 22:7}这预言以更清晰的证言向我们揭示了复活，说：\BibleQuote{我看见死了的人，无论大小，都站在宝座前，案卷就展开了;还有另一卷，即生命册也展开了;死了的人都凭这些案卷所记载的，照他们的行为受了审判。大海交出了其中的死者，地狱也交出了其中的死者。}\BibleRef{默 20:12}那么，我们就不必追问地狱交出、大海归还的那些人是如何复活的。\par

122. 还请听所应许的义人未来的恩宠：\BibleQuote{我听见由宝座那里有一巨大声音说：\InnerQuote{看，天主的帐幕在人间，祂要同人住在一起，他们要作祂的人民，天主亲自要同他们住在一起，作他们的天主。祂要拭去他们眼上的一切泪痕;以后再也没有死亡，再也没有悲伤，没有哀号，没有苦楚，因为先前的都已过去了。}}\BibleRef{默 21:3}\par

123. 现在，你若愿意，请将这生命与那生命对比一下;然后，你若能，就选择在劳苦中、在如此变幻无常的不幸痛苦中、在获得所欲后的餍足中、在伴随快乐而来的厌恶中，使肉身无终止地存在吧。如果天主愿意让这一切永远延续，你会选择它们吗？因为如果生命本身当避开，是为避免烦恼、摆脱痛苦，那么，那紧随其后的是未来复活的永恒欢乐、那里没有过犯的连续、没有犯罪的诱惑的安息，岂不更当追求吗？\par

124. 谁在痛苦中如此忍耐，以至不求死呢？谁在软弱中如此坚忍，以至宁死不欲在衰弱中生活呢？谁在忧愁中如此勇敢，以至不愿借死亡逃避它呢？但是，如果我们在生命持续期间尚且不满，尽管我们知道它有一个期限，那么，要是我们看见肉身的苦难将永远伴随我们，我们对这生命将会变得多么厌倦！因为谁愿意免于死亡呢？或者，什么比一种痛苦的不朽更令人难以忍受呢？\BibleQuote{如果我们只在今生寄望于基督，我们就是众人中最可怜的了。}\BibleRef{格前 15:19}不是因为寄望于基督是可怜的，而是因为基督为那些寄望于祂的人预备了另一种生命。因为今生易犯罪，那生命则是为赏报保留的。\par

124. 我们发现，生命短暂的阶段给我们带来了多少厌倦！男孩渴望成为青年;青年计算着通往更成熟年龄的年数;年轻人，对自己壮年时期的优势不知感恩，却渴望老年的尊荣。于是，所有人自然都生出改变的渴望，因为我们对现在的自己感到不满。最后，就连我们所渴望的事物也令我们厌倦;而我们盼望获得的，一旦获得，便厌恶了。\par

125. 因此，圣善之人并非无故常常哀叹他们在此尘世的漫长居留：\Person{达味}哀叹过，\Person{耶肋米亚}\BibleRef{耶 20:18}哀叹过，\Person{厄里亚}\BibleRef{列上 19:4}也哀叹过。如果我们相信智者，以及那些内有圣神居住的人，他们是在向往更美好的事物；如果我们察考他人的判断，以确定所有人都一致同意，有多少伟人宁愿选择死亡而非忧伤，又有多少伟人宁愿选择死亡而非恐惧！他们确实认为对死亡的恐惧比死亡本身更糟。因此，死亡并非因其自身的邪恶而令人畏惧，而是因生命的苦难，人们反而宁愿选择死亡，因为人们渴望临终者的离去，并回避生者的恐惧。\par

126. 那么，就这样吧。假定复活比此生更值得向往。什么！难道哲学家们自己找到了什么，能让我们比起复活更乐于继续存在吗？事实上，即使是那些声称灵魂不死的人，也不能令我满意，因为他们只允许我得到部分的救赎。如果我不能完全受益，那还能算是何等的恩宠？如果天主的作为在我内消亡，那还能算是何等的生命？如果死亡是自然存在的终点，那么罪人与义人所共有的，又算是何等的正义？所谓真理——灵魂因自身运动且恒动不息而被认为不死——这又算是什么真理？至于那在身体内我们与野兽共有的部分，在身体存在之前究竟发生了什么，或许是不确定的，而真理并不能从这些差异中得出，反而会被摧毁。\par

127. 但那些说我们的灵魂离开这些身体后，会迁移到野兽或其他各种生物的身体中的人，他们的观点难道更可取吗？事实上，哲学家们自己通常也认为，这些不过是诗人的荒诞幻想，就像饮下了\Person{刻耳刻}的药水后可能产生的那样；他们还说，与其说是那些被描绘成遭受了这类事情的人，不如说是那些编造了这类故事的人的理智，仿佛被\Person{刻耳刻}的杯子变成了各种野兽的形态。因为有什么比相信人能被变成野兽的形态更似奇谈呢？然而，统治人的灵魂竟会自取与人如此对立的野兽本性，且本具理性竟能迁入无理性的动物，这岂不比仅身体形态的改变更为怪异吗？你们这些教导这些事的人，摧毁了你们所教导的。因为你们已经放弃了通过魔法咒语来产生这些怪异的转变。\par

128. 诗人戏谑地讲述这些事，哲学家则加以指责，但同时，他们却把那些他们认为对活人而言纯属虚构的事，想象为对死者而言却是真实的。因为那些编造这类故事的人，并非意在断言自己寓言的真相，而是为了嘲笑哲学家的错误：他们认为，那曾习惯于以温和谦卑的意图战胜愤怒的灵魂，如今却能因狮子的狂暴冲动而激愤，按捺不住愤怒，放纵不羁地渴求鲜血、寻求杀戮；或者，那曾像通过君主般的谋略来缓和民众的各种风暴，并以理性的声音使之平静的灵魂，如今却能如狼一般在无路可走的荒野中嚎叫；或者，那曾在重负下呻吟，对着耕犁的劳苦发出悲哀抱怨般低吼的灵魂，如今变成人的模样，在光滑的额头上寻找犄角；又或者，那曾惯于乘快翼翱翔至高天的灵魂，如今却认为飞翔已非自己能力所及，并哀叹在人类身体的重量中变得迟滞。\par

129. 或许你们正是通过某些这样的教导毁掉了\Person{伊卡鲁斯}，因为这位青年在你们的说服下，可能自以为曾是一只鸟。许多老人也因这种手段受骗，甘受剧痛，他们不幸地相信了有关天鹅的寓言，以为自己在以哀歌抚慰痛苦的同时，能够将白发化为茸羽。\par

130. 这些事是多么难以置信！多么可憎！相形之下，按照自然，按照各类果实中所发生的事实去相信，那是多么合宜得多；按照已发生的事件的模式，按照先知的话语，以及基督的天上应许去相信，又是多么合宜！因为有什么比确信天主的工程不会消亡，确信那些按照天主的肖像和模样所造的人，不能被转变为野兽的形态，更美好的呢？因为事实上，并非身体的形态，而是精神的形态，是按照天主的模样造成的。因为，其他种类的受造物既已屈服于人，人又怎能以其自身较优的部分迁移到屈服于自己的动物身上呢？本性不容许这样的事，而即使本性容许，恩宠也不容许。\par

131. 然而，我已看到你们这些外邦人彼此的看法如何，确实，你们这些崇拜野兽的人竟相信自己能被变成野兽，这不应令人感到奇怪。但我更愿你们对适合自己的事有更好的判断，好使你们相信，你们将不会与野兽为伍，而是与天使为伴。\par

灵魂必须离开此生的周遭，摆脱尘世肉身的污浊，奋力奔向那天上的团体。然而，只有圣人才能抵达那里，向天主歌唱赞美（正如我们在先知话中所听到的，那些弹着琴\BibleRef{默 14:2}说：\BibleQuote{上主，全能的天主！你的功行伟大奇妙；万民的君王，你的道路公平正直；谁敢不敬畏你？上主，谁敢不光荣你的名号？因为只有你是圣善的；万民都要前来崇拜你}），\BibleRef{默 15:3} 并瞻仰你的婚宴，主耶稣啊，新娘由尘世被引向天上的事物，万有和谐欢庆，因为\BibleQuote{凡一切血肉都要奔向你}，如今不再受制于易逝之物，却与圣神结合，得见那以细麻、玫瑰、百合、花环装饰的洞房。还有谁的婚事能如此装饰呢？因为这婚事乃是以\OriginalTerm{精修者}{confessors}的紫色绶带、殉道者的鲜血、贞女的百合、司铎的荣冠所装饰的。\par

圣达味最切望的，莫过于能目睹并凝视此事，因为他说：\BibleQuote{我有一事祈求上主，我要恳切请求此事：使我一生的岁月，常居住在上主的殿里，欣赏上主的甘饴慈祥。}\par

信此事是一种愉悦，望此事是一种喜乐；的确，不曾相信它则是一种痛苦，而活于此望中则是恩宠。但若我在这事上错了——宁愿死后与天使为伍，也不愿与禽兽同群——我情愿错下去，只要一息尚存，绝不许人从我骗去这望。\par

因为我除了希望快快到你那里去，我的兄弟，希望你的离去不会在我们之间造成长久的分离，并希望藉着你的代祷，你能快快召叫我——这切望你的人——之外，还有什么安慰存留呢？因为谁不当超乎一切地切愿自己\BibleQuote{这必朽坏的穿上那不杇坏的，这必死的穿上那不死的}\BibleRef{格前 15:53}呢？这样，我们这因肉身的软弱而屈服于死亡的人，一旦被提升超越本性，就不再畏惧死亡了。\par

\SourceRecord{Translated by H. de Romestin, E. de Romestin and H.T.F. Duckworth.；Nicene and Post-Nicene Fathers, Second Series；Vol. 10.；Edited by Philip Schaff and Henry Wace.；Buffalo, NY: Christian Literature Publishing Co.,；1896.；<http://www.newadvent.org/fathers/34032.htm>.}
